\documentclass[11pt,reqno]{amsart}
\usepackage[T1]{fontenc}
\usepackage{lmodern}
\usepackage[a4paper,margin=3cm]{geometry}
\usepackage{amsmath,amssymb,mathtools,amsthm}
\usepackage{mathrsfs}
\usepackage{needspace}
\usepackage{microtype}
\usepackage{enumitem}
\usepackage{booktabs}
\usepackage[colorlinks=true,linkcolor=blue,citecolor=blue,urlcolor=blue,hypertexnames=false]{hyperref}
\usepackage{fancyhdr}
\numberwithin{equation}{section}
\setcounter{tocdepth}{1}
\newtheorem{theorem}{Theorem}[section]
\newtheorem{lemma}[theorem]{Lemma}
\newtheorem{proposition}[theorem]{Proposition}
\newtheorem{corollary}[theorem]{Corollary}
\theoremstyle{definition}
\newtheorem{definition}[theorem]{Definition}
\theoremstyle{remark}
\newtheorem{remark}[theorem]{Remark}
\newtheorem{example}[theorem]{Example}
\newcommand{\Z}{\mathbb Z}
\newcommand{\R}{\mathbb R}
\newcommand{\T}{\mathbb T}
\newcommand{\N}{\mathbb N}
\newcommand{\C}{\mathbb C}
\newcommand{\E}{\mathbb E}
\newcommand{\La}{\mathcal L}
\newcommand{\one}{\mathbf 1}
\newcommand{\norm}[1]{\lVert#1\rVert}
\newcommand{\abs}[1]{\lvert#1\rvert}
\newcommand{\ip}[2]{\langle#1,#2\rangle}
\DeclareMathOperator{\supp}{supp}
\DeclareMathOperator{\Rea}{Re}
\DeclareMathOperator{\Ima}{Im}
\DeclareMathOperator{\Maj}{Maj}
\newcommand{\eps}{\varepsilon}
\newcommand{\cMPS}{c_{\mathrm{MPS}}}
\newcommand{\doi}[1]{\href{https://doi.org/#1}{\texttt{doi:#1}}}
\newcommand{\arxiv}[1]{\href{https://arxiv.org/abs/#1}{\texttt{arXiv:#1}}}
\setlength{\parskip}{0em}
\setlist{itemsep=0pt,topsep=0.3em,parsep=0pt,partopsep=0pt}
\setlist[enumerate]{leftmargin=*,label=\textup{(\roman*)}}
\allowdisplaybreaks[1]
\pagestyle{fancy}
\fancyhf{}
\renewcommand{\headrulewidth}{0pt}
\setlength{\headheight}{13pt}
\setlength{\footskip}{13pt}
\fancyhead[C]{\ifodd\value{page}{\small XIYU HU}\else{\small LITTLEWOOD CONSTANTS AND STRUCTURAL INVERSE THEOREMS}\fi}
\fancyfoot[C]{\thepage}
\setlength{\emergencystretch}{2em}
\hypersetup{pdftitle={An improved Littlewood constant and structural inverse theorems},pdfauthor={Xiyu Hu},pdfsubject={Fourier L1 norms, additive energy, and structured integer sets},pdfkeywords={Littlewood conjecture, inverse theorem, additive energy, outer functions, Sidon sets}}
\title[Littlewood constants and structural inverse theorems]{An improved Littlewood constant and structural inverse theorems}
\author{Xiyu Hu}
\address{University of Chinese Academy of Sciences}
\email{hxypqr@gmail.com}
\date{}
\subjclass[2020]{Primary 42A05, 11B30; Secondary 42A55, 43A46}
\keywords{Littlewood conjecture, Fourier $L^1$ norm, additive energy, outer functions, Sidon sets, Freiman isomorphisms}
\begin{document}
\begin{abstract}
For a finite set $A\subset\Z$, write
$\La(A)=\int_0^1|\sum_{a\in A}e^{2\pi iat}|\,dt$.
We prove the uniform lower bound
$\La(A)\ge(c_*-o(1))\log|A|$, where $c_*>0.2459209$ is explicit.
We also establish an almost-covering inverse theorem for an explicit class
of integer sets generated from dense separated blocks and long arithmetic
block inflations by cyclic Sidon-fibre decompositions of arbitrary depth.
If $A$ belongs to this class and $\La(A)\le K\log|A|$, then, for every
$\eps>0$, all but $\eps|A|$ points are covered by boundedly many disjoint,
linear-sized subsets of bounded doubling. The constants depend only on
$K$, $\eps$, and the fixed terminal-geometry parameters, and are independent
of the decomposition depth, branching, moduli, and diameter.
The same conclusion holds under high-order Freiman modelling.
\end{abstract}
\maketitle
\tableofcontents

\Needspace{7\baselineskip}
\section{Introduction}
\label{sec:introduction}

Let $e(t)=\exp(2\pi it)$ and, for a finite set of integers $A$, put
\begin{equation}\label{eq:norm-definition}
 F_A(t)=\sum_{a\in A}e(at),\qquad
 \La(A)=\int_{\T}|F_A(t)|\,dt,\qquad \T=\R/\Z.
\end{equation}
The measure on $\T$ is the probability Haar measure, and all logarithms in
this paper are natural. Littlewood's conjecture, proved independently by
Konyagin~\cite{Konyagin} and McGehee, Pigno, and Smith~\cite{MPS}, states that
$\La(A)\gg\log|A|$. The order of magnitude is attained by arithmetic
progressions. Indeed, if $P$ is a progression of length $N$, then
\begin{equation}\label{eq:ap-asymptotic}
 \La(P)=\frac4{\pi^2}\log N+O(1).
\end{equation}
Translation and nonzero integer dilation preserve $\La$, so this expression
does not depend on the location or common difference of $P$.

There are two natural questions beyond the logarithmic lower bound. The
first is quantitative: can its leading constant be increased towards
$4/\pi^2$? The second is structural: what can a set look like when its
Fourier $L^1$ norm is only a constant multiple of $\log N$, rather than a
quantity comparable to the elementary upper bound $\sqrt N$? These
questions are related through duality, but a quantitative improvement to
a bounded test function does not by itself provide a structural inverse
theorem. This paper treats both questions, and keeps their conclusions
and hypotheses separate.

\subsection{Previous work and the connection with Bloom--Green}

The stronger assertion that an arithmetic progression minimizes $\La(A)$
among all $N$-element sets, as well as its leading-constant version, is
usually called the strong Littlewood conjecture; see Hardy--Littlewood
\cite{HardyLittlewood} and Zygmund~\cite{Zygmund}. The MPS proof gives a
weighted inequality for arbitrary complex coefficients, not merely the
indicator-set case. Stegeman~\cite{Stegeman} and Yabuta~\cite{Yabuta}
refined the numerical constant in that argument, with Yabuta obtaining a
coefficient approximately $0.12959$. Fournier~\cite{Fournier} discusses the
bounded-test-function methods underlying the classical proofs.

Bloom and Green~\cite{BG} revisit the MPS construction with two purposes.
They obtain the coefficient $0.170934\ldots$, and, by examining the
failure of a putative larger lower bound, prove the following inverse
statement. If $|A|=N$ and $\La(A)\le K\log N$, then for
$0<\delta\le1/2$ there is an initial segment $A'\subseteq A$ with
\begin{equation}\label{eq:BG-baseline}
 |A'|\gg N^{1-\delta},\qquad
 E(A')\gg(\delta/K)^2|A'|^3,
\end{equation}
where
\[
 E(B)=\#\{(b_1,b_2,b_3,b_4)\in B^4:b_1+b_2=b_3+b_4\}
     =\int_{\T}|F_B|^4
\]
is the additive energy. Their analysis exposes an energy square root in
the error of the ordered MPS recursion. Standard additive-combinatorial
results then turn \eqref{eq:BG-baseline} into small-doubling subsets and
arithmetic-progression consequences. Their result is unrestricted: it
applies to every finite integer set.

Earlier structural work addressed rather different descriptions of
complexity. Petridis~\cite{Petridis} and Hanson~\cite{Hanson} show that
sufficiently genuine multidimensional structure increases the Fourier
norm. Pichorides~\cite{Pichorides}, and independently Bedert~\cite{Bedert},
relate the norm to the size of dissociated subsets. In another direction,
the finite-group theorem of Green and Sanders~\cite{GreenSanders}
classifies integer-valued functions of bounded Fourier algebra norm by
bounded combinations of coset indicators. The latter has a constant
spectral budget, unlike the logarithmic budget of the integer
Littlewood problem. Bourgain~\cite{Bourgain} and Bedert~\cite{Bedert}
also connect these questions to large sum-free subsets.

Our starting point is the same dual principle as in Bloom--Green: a
bounded test function tries to accumulate contributions, and the
obstruction to accumulation contains information. For the numerical
problem we retain the entire product of the damping multipliers. For the
inverse problem we arrange that an obstruction is inherited by every
remaining family of blocks, while all pairings continue to use the
\emph{original} target. This distinction allows repeated structural
extraction without assuming that deleting frequencies decreases an
$L^1$ norm.

\subsection{The quantitative theorem}

For $p>1$ define
\begin{equation}\label{eq:sinc-intro}
 I_p=\int_{\R}\left|\frac{\sin\pi x}{\pi x}\right|^p\,dx,
 \qquad c_{2r}=I_{2r},\qquad
 c_{2r+1}=\sqrt{I_{2r}I_{2r+2}}\quad(r\ge1).
\end{equation}
Put $a=9/10$ and
\begin{align}
 \mathcal E_a&=\frac{55a^2}{108}
 +\frac19\left\{
 \frac23\left(\frac{5a^4}{6}-\frac{2a^3}{3}\right)
 +2\sum_{k=5}^{\infty}\frac{c_ka^k}{k}\right\},
 \label{eq:E-intro}\\
 c_*&=\frac{9}{10\log9}\left(1-\sqrt{\mathcal E_{9/10}/3}\right).
 \label{eq:cstar}
\end{align}
The series is absolutely convergent, since $0<c_k\le1$.

\begin{theorem}[An explicit Littlewood coefficient]\label{thm:constant}
For every finite set $A\subset\Z$ of cardinality $N$,
\begin{equation}\label{eq:constant-result}
 \La(A)\ge(c_*-o(1))\log N,
 \qquad c_*>0.2459209.
\end{equation}
The $o(1)$ is uniform over all $N$-element integer sets.
\end{theorem}

The parameter choices in \eqref{eq:E-intro} are explicit rather than
numerically optimized. The proof has three ingredients. First, an outer
multiplier of modulus $1-a|F_B|/|B|$ uses exactly the pointwise room left
by the source term $aF_B/|B|$. Second, the error of a product is estimated
before its factors are separated; its second-order expansion retains the
loss of squared modulus. Third, the same scalar variable controls both
terms, and interval moment bounds handle the only tail level at which a
nonconstant majorant is needed. We provide a self-contained proof of the
moment comparison using symmetric chains, and an exact-rational
certificate for the displayed decimal.

\subsection{The structural theorem and its precise class}

The conclusion of our inverse theorem is an almost-covering by large
small-doubling sets. Its hypotheses specify how the original integer set
is assembled; they are not deduced from small $\La$. We give the complete
definition here so that the scope of the main statement is visible from
the introduction.

\begin{definition}[Additive Sidon sets]\label{def:sidon}
A subset $R$ of an abelian group is \emph{Sidon} if
$r_1+r_2=r_3+r_4$ with $r_i\in R$ implies equality of the two unordered
pairs, with repetitions allowed. This is the additive $B_2$ condition,
not the harmonic-analytic definition of a Sidon spectrum.
\end{definition}

Fix $0<\beta\le1$ and $\Delta\ge1$. We use two terminal geometries,
together with singletons. Their affine copies are also allowed.

\begin{definition}[Terminal sets]\label{def:terminals}
A nonempty finite integer set is a terminal set if it is a singleton,
or an affine copy under $x\mapsto u+vx$, $v\in\Z\setminus\{0\}$, of one
of the following sets.
\begin{enumerate}
\item A \emph{dense separated block set} $C$, of size $n$, has an exact
representation
\begin{equation}\label{eq:terminal-packet}
 C=\bigcup_{i=1}^R(b_i+C_i),\qquad C_i\subseteq[-D,D]\cap\Z,
\end{equation}
where $D,M$ are positive integers and
\begin{equation}\label{eq:terminal-packet-conditions}
 |C_i|\ge M\ge n^\beta,\qquad D\le\Delta M,
 \qquad |b_i-b_j|>4D\quad(i\ne j).
\end{equation}
The centres need not lie in a common progression and the $C_i$ need not
be intervals.
\item A \emph{long arithmetic block inflation} is
\begin{equation}\label{eq:terminal-inflation}
 A_M(\mathcal C)=
 \{b+q(Mu+r):(b,u)\in\mathcal C,\ 0\le r<M\},
\end{equation}
where $q,M\ge1$, $\varnothing\ne\mathcal C\subset(\Z/q\Z)\times\Z$ is
finite, $b$ is represented in $\{0,\ldots,q-1\}$, and
$M\ge V:=|\mathcal C|$. Its cardinality is $MV$.
\end{enumerate}
\end{definition}

\begin{definition}[The assembly class]\label{def:class}
Let $\mathfrak G_{\beta,\Delta}$ consist of the nonempty finite integer
sets represented by finite rooted trees with terminal sets as leaves.
At every internal vertex, after an affine normalization, its set is the
exact union
\begin{equation}\label{eq:node}
 A_v=\bigcup_{r\in R_v}(r+q_v A_{v,r}),\qquad q_v\ge2,
\end{equation}
where $R_v\subset\Z/q_v\Z$ is the \emph{entire} nonempty residue support,
is Sidon, and the children are the nonempty sets $A_{v,r}$. The images
of the children partition the parent. One-child vertices are allowed.
There is no bound on depth, branching numbers, moduli, integer diameter,
or the number of leaves.
\end{definition}

\begin{theorem}[An almost-covering inverse theorem]\label{thm:inverse}
Fix $K\ge1$, $0<\eps<1/2$, $0<\beta\le1$, and $\Delta\ge1$. There are
$N_0$, $c>0$, and $C<\infty$, depending only on these parameters, with the
following property. If
\[
 A\in\mathfrak G_{\beta,\Delta},\qquad |A|=N\ge N_0,
 \qquad \La(A)\le K\log N,
\]
then there are pairwise disjoint nonempty subsets $B_1,\ldots,B_J\subseteq A$
such that
\begin{equation}\label{eq:cover-conclusion}
 \left|A\setminus\bigcup_{j=1}^J B_j\right|\le\eps N,
 \qquad J\le C,\qquad |B_j|\ge cN,\qquad
 |B_j+B_j|\le C|B_j|.
\end{equation}
The constants do not depend on any of the tree parameters excluded in
Definition~\ref{def:class}.
The same conclusion holds if an order-$N$ Freiman-isomorphic image of
$A$ belongs to $\mathfrak G_{\beta,\Delta}$.
\end{theorem}

Here a Freiman isomorphism of order $h$ is a bijection preserving in both
directions all equalities between sums of at most $h$ terms on each side,
with equal numbers of terms and with repetitions. The last clause is
justified by a uniform estimate for the Fourier norms of \emph{all}
corresponding subsets; no order-two invariance of the norm is assumed.

\begin{corollary}[Macroscopic hereditary energy]\label{cor:main-hered}
Fix $K,\beta,\Delta$ and $0<\delta\le1$. Under the geometric or modelling
hypothesis of Theorem~\ref{thm:inverse}, if $N$ is sufficiently large and
$\La(A)\le K\log N$, every $Y\subseteq A$ with $|Y|\ge\delta N$ satisfies
\[
 E(Y)\ge\rho_{K,\beta,\Delta,\delta}|Y|^3
\]
for a positive constant independent of $A,N,Y$.
\end{corollary}

Theorem~\ref{thm:inverse} has a stronger conclusion than
\eqref{eq:BG-baseline} on the specified class, but has a stronger
hypothesis as well. It is not an unrestricted improvement of the
Bloom--Green inverse theorem. In particular, we do not prove that a
small-norm integer set possesses the tree in Definition~\ref{def:class},
even after high-order modelling. The pieces in \eqref{eq:cover-conclusion}
are small-doubling sets, not necessarily single progressions, and their
number can depend on $\eps$. Thus the result also stops short of the
signed-progression conjecture of Petridis and Bloom--Green.

\subsection{How the two arguments fit together}

In both parts of the paper, the main issue is to keep the information
that a norm estimate would otherwise discard. In the quantitative part,
separately maximizing the outer-function error and the surviving modulus
loses their dependence on the same variable. The whole-product estimate
restores that dependence. In the structural part, replacing all failed
amplifications by a single energy estimate loses the possibility of
iteration. We instead prove an ambient-target hereditary inequality
\begin{equation}\label{eq:hereditary-intro}
 E(S)\ge\eta\frac{|S|^4}{|\Lambda|}
 \qquad(\varnothing\ne S\subseteq\Lambda)
\end{equation}
for the integer labels of separated blocks. Bounded majority polynomials
force a positive density of mixed correlations, and finite Fourier
bandwidth turns those correlations into relations with a bounded carry.
The original target is never replaced by the remaining block union.

Long arithmetic blocks provide a different route to the same invariant.
Taking a discrete boundary removes the long direction, and the coefficient
of its logarithmic contribution is a constant-scale spectral norm.
Integer-valuedness then bounds the number of common boundary heights.
Finite-group structure can be lifted back to integers while retaining
the possible carries. Both local arguments feed into a hereditary
energy-to-covering principle based on BSG.

A final obstruction to iteration is the multiplication of losses along
an arbitrarily deep tree. We avoid it through the coefficient inequality
\[
 \left\|\sum_{r\in R}a_r\chi_r\right\|_1
 \ge\left(\sum_{r\in R}|a_r|^{20}\right)^{1/20}
 \quad\text{for additive Sidon }R.
\]
The exponent is not optimized; the constant one is essential. It gives
an exact budget for all leaves at once, even for very unbalanced trees.
Only boundedly many macroscopic leaves survive the logarithmic norm
bound, and the two local covering theorems complete the proof.

Section~\ref{sec:preliminaries} records the external inputs.
Sections~\ref{sec:moments}--\ref{sec:constant-proof} prove
Theorem~\ref{thm:constant}. Sections~\ref{sec:hered}--\ref{sec:gluing}
prove Theorem~\ref{thm:inverse}; Section~\ref{sec:transport} treats modelling
and small exceptional perturbations. Examples and the remaining general
problem are discussed in Section~\ref{sec:scope}. The appendices give
an exact numerical certificate and a one-sided obstruction, together
with a two-sided interval model for the sharp constant.

\Needspace{7\baselineskip}
\section{Notation and standard inputs}\label{sec:preliminaries}

For a function $f$ on $\T$, its Fourier coefficient is
$\widehat f(n)=\int_\T f(t)e(-nt)\,dt$. For finitely supported functions on
an abelian group, convolution has its usual counting-measure convention.
The transform $F_X$ of an indicator is the sum of the associated
characters on the compact dual. All dual Haar measures, including those
on finite groups, have total mass one. Pairings are linear in the first
entry: $\ip fg=\int f\overline g$. The normalized energy is
$\omega(X)=E(X)/|X|^3$. We use $c,C$ for positive absolute constants that
may change, with subscripts indicating additional dependence.

The identity $\int|F_X|^2=|X|$ and Fourier inversion imply
\begin{equation}\label{eq:trivial-norms}
 1\le\La(X)\le\sqrt{|X|}\quad(X\ne\varnothing).
\end{equation}
Also, by Cauchy--Schwarz applied to the representation function of $X+X$,
\begin{equation}\label{eq:energy-doubling}
 E(X)\ge\frac{|X|^4}{|X+X|}.
\end{equation}
These statements apply to finite subsets of any discrete abelian group.

Apart from standard Fourier and Hardy-space facts, we use the following
two established results as external inputs. The further analytic and
combinatorial implications needed for our main statements are proved below.

\begin{theorem}[McGehee--Pigno--Smith~\cite{MPS}]\label{thm:MPS}
There is an absolute $\cMPS>0$ such that, whenever
$n_1<\cdots<n_m$ are integers and $u_1,\ldots,u_m\in\C$,
\begin{equation}\label{eq:MPS-weighted}
 \int_\T\left|\sum_{j=1}^m u_je(n_jt)\right|\,dt
 \ge\cMPS\sum_{j=1}^m\frac{|u_j|}{j}.
\end{equation}
In particular, after decreasing $\cMPS$ if necessary,
$\La(X)\ge\cMPS\log(|X|+1)$ for every nonempty finite integer set.
\end{theorem}
The complex-coefficient formulation \eqref{eq:MPS-weighted} is also
stated explicitly in Hanson~\cite[Theorem 1.1]{Hanson}. We will need it
with some coefficients equal to zero, which is harmless: removing them
only decreases their indices on the right.

\begin{theorem}[Polynomial Balog--Szemer\'edi--Gowers]\label{thm:BSG}
There are absolute $c,C>0$ such that, for a finite nonempty subset $Y$ of
an abelian group and $0<\gamma\le1$, the inequality
$E(Y)\ge\gamma|Y|^3$ yields a subset $Z\subseteq Y$ with
\[
 |Z|\ge c\gamma^C|Y|,
 \qquad |Z+Z|\le C\gamma^{-C}|Z|.
\]
\end{theorem}
We use this standard qualitative-polynomial form; see
Tao--Vu~\cite[Theorem 2.31]{TaoVu} and the quantitative refinements of
Reiher--Schoen~\cite{ReiherSchoen}. The latter are formulated with a
small difference set. The standard sum--difference inequalities convert
that conclusion to the displayed form, with a change in the absolute
powers. No optimization of BSG exponents is used here.

For $q\ge2$ and $r\in\{0,\ldots,q-1\}$ set
$A_r=\{m:r+qm\in A\}$. Character orthogonality gives
\begin{equation}\label{eq:residue-projection}
 \frac1q\sum_{j=0}^{q-1}e(-rj/q)F_A(t+j/q)
   =e(rt)F_{A_r}(qt).
\end{equation}
Consequently $\La(A_r)\le\La(A)$. Moreover
\begin{equation}\label{eq:affine-invariance}
 \La(u+vA)=\La(A)\qquad(u,v\in\Z,\ v\ne0).
\end{equation}
Both identities follow from Haar invariance, including for negative
$v$. The contraction in \eqref{eq:residue-projection} does not control
$|A_r|$, and will not be used to discard a small fibre without a separate
mass estimate.

\Needspace{7\baselineskip}
\section{Interval bounds for all even moments}\label{sec:moments}

The numerical argument needs precise moment constants. Although these
bounds are related to classical Fourier rearrangement inequalities
\cite{Gabriel,HardyLittlewood}, for indicator sets they follow from an
elementary antichain argument. We give the argument to make the precise
extremal input, including its constant, independent of a rearrangement
convention.

\begin{lemma}[Symmetric chains]\label{lem:chains}
The product $\{0,\ldots,n-1\}^d$, partially ordered coordinatewise, has a
partition into saturated chains symmetric about rank $d(n-1)/2$.
In particular every antichain has size at most the largest rank level.
\end{lemma}
\begin{proof}
A single chain has the required property. To pass from a symmetric chain
decomposition of a ranked poset to one of its product with a chain, it
suffices to decompose a rectangle $\{0,\ldots,p\}\times\{0,\ldots,q\}$.
Assume $p\ge q$, interchanging the coordinates otherwise. For
$0\le j\le q$, take the chain
\[
 \{(i,j):0\le i\le p-j\}
 \ \cup\ \{(p-j,k):j<k\le q\}.
\]
These chains are disjoint and cover the rectangle. The first and last
ranks of the $j$th chain are $j$ and $p+q-j$. Thus each is symmetric.
If the original chain starts at an offset rank, the same construction
preserves symmetry about the rank of the full product. Induction proves
the decomposition. Every symmetric chain meets a middle rank exactly
once, and every antichain meets a chain at most once. This proves the
last assertion. This is the symmetric-chain construction of de Bruijn,
van Ebbenhorst Tengbergen, and Kruyswijk~\cite{BTK}.
\end{proof}

\begin{proposition}[Even-moment comparison]\label{prop:moments}
If $B\subset\Z$ has $n$ elements and $r\ge1$, then
\begin{equation}\label{eq:even-max}
 \int_\T|F_B|^{2r}\le
 [z^{r(n-1)}](1+z+\cdots+z^{n-1})^{2r}
 =\int_\T|F_{\{0,\ldots,n-1\}}|^{2r}.
\end{equation}
In particular
\begin{equation}\label{eq:energy-max}
 E(B)\le\frac{2n^3+n}{3}.
\end{equation}
\end{proposition}
\begin{proof}
Write $B=\{b_0<\cdots<b_{n-1}\}$. An additive relation counted by the
left side of \eqref{eq:even-max} is represented by a point
$x\in\{0,\ldots,n-1\}^{2r}$ satisfying
\[
 b_{x_1}+\cdots+b_{x_r}
 -b_{n-1-x_{r+1}}-\cdots-b_{n-1-x_{2r}}=0.
\]
The expression on the left strictly increases when any coordinate
increases. Its zero set is therefore an antichain. By
Lemma~\ref{lem:chains}, it has size at most the central rank of the grid,
which is the coefficient in \eqref{eq:even-max}. For the interval, every
point of that rank satisfies the required equality, so equality holds.
For $r=2$, the interval difference representation function is
$n-|h|$ for $|h|<n$. Its squared sum is
$n^2+2\sum_{j=1}^{n-1}j^2=(2n^3+n)/3$.
\end{proof}

\begin{corollary}[Uniform normalized moments]\label{cor:normalized-moments}
Let $u_B=|F_B|/n$ and let $c_k$ be as in \eqref{eq:sinc-intro}. For each
fixed integer $k\ge2$,
\begin{equation}\label{eq:normalized-moments}
 n\int_\T u_B^k\le c_k+o_k(1)\quad(n\longrightarrow\infty),
\end{equation}
uniformly in $B$. For every $k\ge2$ and every $B$, the left side is at
most one. For $r\ge1$,
\begin{equation}\label{eq:even-sinc-formula}
 I_{2r}=\frac1{(2r-1)!}
 \sum_{j=0}^r(-1)^j\binom{2r}{j}(r-j)^{2r-1}.
\end{equation}
\end{corollary}
\begin{proof}
For the interval polynomial $D_n(t)=\sum_{j=0}^{n-1}e(jt)$, change variables
$x=nt$ on $[-1/2,1/2]$. The rescaled integrand is
\[
 \left|\frac{\sin\pi x}{n\sin(\pi x/n)}\right|^{2r}
 \one_{[-n/2,n/2]}(x).
\]
It converges pointwise to $|\sin(\pi x)/(\pi x)|^{2r}$ and is bounded
by a constant times $\min(1,|x|^{-2r})$. Dominated convergence and
Proposition~\ref{prop:moments} prove \eqref{eq:normalized-moments} for even
$k$. Cauchy--Schwarz between the adjacent even moments gives the odd
case. The universal bound follows from $0\le u_B\le1$ and
$n\int u_B^2=1$.

For completeness, the coefficient on the right of \eqref{eq:even-max}
equals
\[
 \sum_{j=0}^{r-1}(-1)^j\binom{2r}{j}
 \binom{(r-j)n+r-1}{2r-1},
\]
where binomial coefficients with too small a top entry are interpreted
as zero. Divide by $n^{2r-1}$ and let $n\to\infty$. Each term has the
limit $(r-j)^{2r-1}/(2r-1)!$. The additional $j=r$ term in
\eqref{eq:even-sinc-formula} vanishes. This proves the formula and also
shows directly that the limiting moment constants are rational.
\end{proof}

\Needspace{7\baselineskip}
\section{Outer multipliers and the improved constant}\label{sec:constant-proof}

\subsection{An exact pointwise budget}

We work with multipliers whose Fourier support is nonpositive. This
orientation is convenient for increasing initial segments. We use standard
Hardy-space boundary convergence as in Duren~\cite{Duren}. The following
construction is the boundary-value form of an outer function, reflected
from the usual analytic convention.

\begin{lemma}[Outer damping]\label{lem:outer}
Let $0<a<1$, and let $u:\T\to[0,1]$ be measurable. There is an
$L^\infty$ function $M$ with
\begin{equation}\label{eq:outer-data}
 \supp\widehat M\subseteq\Z_{\le0},\qquad
 |M|=1-au\ \text{a.e.},\qquad
 \widehat M(0)=\exp\!\left(\int_\T\log(1-au)\right)>0.
\end{equation}
If $M_1,\ldots,M_v$ are such multipliers and $P=\prod_jM_j$, then
\begin{equation}\label{eq:outer-product-mean}
 \int_\T P=\prod_j\widehat M_j(0).
\end{equation}
\end{lemma}
\begin{proof}
Put $\ell=\log(1-au)$, a real bounded function. For $0<\rho<1$ define
\[
 M_\rho(t)=\exp\!\left(\widehat\ell(0)
       +2\sum_{k<0}\widehat\ell(k)\rho^{|k|}e(kt)\right).
\]
The series converges absolutely for each $\rho$. Its real part before
exponentiation is the Poisson integral of $\ell$, so
$1-a\le |M_\rho|\le1$. These are the radial values of a bounded analytic
function in the reflected disc. The radial boundary function $M$ exists
almost everywhere and in every finite $L^p$. Poisson convergence gives
$|M|=e^\ell=1-au$, while its constant coefficient is
$e^{\widehat\ell(0)}$. This proves \eqref{eq:outer-data}. Products are
again bounded reflected-analytic functions, and their values at the
origin multiply. Equivalently, radial approximation followed by
$L^1$ convergence proves \eqref{eq:outer-product-mean}.
\end{proof}

For a finite nonempty $B$, let $M_{a,B}$ denote this multiplier for
$u_B=|F_B|/|B|$. Whenever $|G|\le1$, the update
\begin{equation}\label{eq:outer-update}
 G\longmapsto GM_{a,B}+aF_B/|B|
\end{equation}
preserves the bound one, because the two summands have total modulus at
most $(1-au_B)+au_B=1$. Unlike an exponential majorant for the modulus,
this update leaves no unused pointwise room.

\subsection{Ordered pairing and the whole-product error}

Write $A=\{a_1<\cdots<a_N\}$ and let $A_i$ be increasing initial
segments of sizes $n_i$. For now the number of segments is arbitrary.
Set $M_i=M_{a,A_i}$ and
\begin{equation}\label{eq:test-function}
 P_i=\prod_{j>i}M_j,\qquad
 G=a\sum_i\frac{F_{A_i}}{n_i}P_i.
\end{equation}
Iterating \eqref{eq:outer-update} from zero proves $\norm G_\infty\le1$.
The frequency orientation gives the following exact localization of the
pairing.

\begin{lemma}[The initial-segment pairing]\label{lem:prefix-pairing}
For every initial segment $A_i$ and every $P$ with nonpositive Fourier
support,
\begin{equation}\label{eq:prefix-error}
 \left|\ip{F_A}{(F_{A_i}/n_i)P}-1\right|
 \le
 \left(\frac{E(A_i)+n_i^2}{2n_i^2}\right)^{1/2}
 \norm{P-1}_2.
\end{equation}
Here $P$ may be any $L^2$ function for which the displayed pairing is
formed; in particular it may be the bounded product $P_i$.
\end{lemma}
\begin{proof}
Let $p_s=\widehat P(s)$. In the pairing, a coefficient $p_s$ can occur
only when $a=b+s$, with $a\in A$, $b\in A_i$, and $s\le0$. Then
$a\le b$, so the initial-segment property implies $a\in A_i$.
Consequently the pairing is unchanged if $F_A$ is replaced by $F_{A_i}$.
For $s\in\Z$ write
\[
 d_s=\frac1{n_i}\#\{(a,b)\in A_i^2:a-b=s\}.
\]
Then $d_0=1$, $d_s=d_{-s}$, and
$\sum_sd_s^2=E(A_i)/n_i^2$. Subtracting the contribution of the constant
function $P=1$, Cauchy--Schwarz on $s\le0$ gives
\[
 \left|\sum_{s\le0}d_s\overline{(p_s-\one_{s=0})}\right|
 \le\left(\sum_{s\le0}d_s^2\right)^{1/2}\norm{P-1}_2.
\]
The sum of squares in the last expression is
$(E(A_i)/n_i^2+1)/2$, as claimed. All sums involving $d_s$ are finite,
so no issue of pointwise Fourier convergence arises.
\end{proof}

Define, for $0\le x\le1$,
\begin{equation}\label{eq:h-q}
 h_a(x)=a^2x^2-2ax-2\log(1-ax),\qquad
 q_a(x)=2ax-a^2x^2.
\end{equation}
The function $h_a$ has the positive power series
\begin{equation}\label{eq:h-expansion}
 h_a(x)=2a^2x^2+2\sum_{k\ge3}\frac{a^kx^k}{k}.
\end{equation}

\begin{lemma}[Whole-product estimate]\label{lem:whole-product}
Let $M_j$ be the multipliers in Lemma~\ref{lem:outer}, associated to
$u_j:\T\to[0,1]$, and put $P=\prod_jM_j$. Then
\begin{equation}\label{eq:whole-product}
 \norm{P-1}_2^2
 \le\sum_j\int_\T h_a(u_j)
       +\sum_{j<k}\int_\T q_a(u_j)q_a(u_k).
\end{equation}
\end{lemma}
\begin{proof}
Set
\[
 W=\prod_j(1-au_j),\qquad
 S=\sum_j\int_\T-\log(1-au_j).
\]
By \eqref{eq:outer-product-mean}, $\int P=e^{-S}$, a positive real
number. Hence
\begin{equation}\label{eq:product-identity}
 \norm{P-1}_2^2=\int_\T W^2+1-2e^{-S}.
\end{equation}
We have $W^2=\prod_j(1-q_a(u_j))$ and $0\le q_a(u_j)\le1$. The second
Bonferroni inequality gives the pointwise upper bound
\[
 W^2\le1-\sum_jq_a(u_j)+\sum_{j<k}q_a(u_j)q_a(u_k).
\]
For instance, this is the usual two-term bound for the probability that
none of independent events of probabilities $q_a(u_j)$ occurs, applied
at a fixed point of $\T$. Since $1-e^{-S}\le S$, substitution in
\eqref{eq:product-identity} proves \eqref{eq:whole-product}, using
$-q_a(x)-2\log(1-ax)=h_a(x)$.
\end{proof}

The point of \eqref{eq:whole-product} is not merely the retention of a
cross term. The quantities $h_a(u)$ and $q_a(u)^2$ depend on the same
variable. They cannot in general be simultaneously as large as their
separate uniform upper bounds.

\subsection{A coupled scalar majorant}

We now fix $a=9/10$ and impose the exact scale ratio
$n_{i+r}=9^rn_i$. For a fixed $i$, let
\[
 D_r=n_{i+r}\int_\T h_a(u_{A_{i+r}}),\qquad
 Q_r=n_{i+r}\int_\T q_a(u_{A_{i+r}})^2,
 \qquad w_r=3^{-r}.
\]
Cauchy--Schwarz on each cross term of \eqref{eq:whole-product} gives
\begin{equation}\label{eq:weighted-tail}
 n_i\norm{P_i-1}_2^2
 \le\sum_rw_r^2D_r+\sum_{r<s}w_rw_s\sqrt{Q_rQ_s}.
\end{equation}
All indices here range only over the existing tail of the chain.
Choose the infinite sequence
\begin{equation}\label{eq:v-weights}
 v_1=\frac14,\qquad v_r=3^{-r}\quad(r\ge2),\qquad
 V_0=\sum_{r\ge1}v_r=\frac5{12},\qquad
 b_r=\frac12\left(\frac{V_0}{v_r}-1\right).
\end{equation}
In particular $b_1=1/3$ and $b_r\ge b_2=11/8$ for $r\ge2$.
Weighted Cauchy--Schwarz, with the possibly unused $v_r$ retained in
$V_0$, gives
\[
 \left(\sum_rw_r\sqrt{Q_r}\right)^2
 \le V_0\sum_r\frac{w_r^2Q_r}{v_r}.
\]
It follows from \eqref{eq:weighted-tail} that
\begin{equation}\label{eq:coupled-tail}
 n_i\norm{P_i-1}_2^2
 \le\sum_rw_r^2(D_r+b_rQ_r).
\end{equation}

\begin{lemma}[Scalar coupling]\label{lem:scalar}
Let $a=9/10$. For a nonempty finite integer set $B$, $n=|B|$, put
$u=|F_B|/n$. If $b\ge11/8$, then
\begin{equation}\label{eq:scalar-large-b}
 n\int_\T\bigl(h_a(u)+bq_a(u)^2\bigr)\le(2+4b)a^2.
\end{equation}
For $b=1/3$, uniformly as $n\to\infty$,
\begin{equation}\label{eq:scalar-first}
 n\int_\T\bigl(h_a(u)+\tfrac13q_a(u)^2\bigr)
 \le\frac{10a^2}{3}+\Delta_a+o(1),
\end{equation}
where
\begin{equation}\label{eq:Delta}
 \Delta_a=\frac23\left(\frac{5a^4}{6}-\frac{2a^3}{3}\right)
                 +2\sum_{k\ge5}\frac{c_ka^k}{k}.
\end{equation}
\end{lemma}
\begin{proof}
By Parseval, $d\nu=nu^2\,dt$ is a probability measure. Define
\[
 \Phi_b(x)=\frac{h_a(x)}{x^2}+b(2a-a^2x)^2\quad(x>0),\qquad
 \Phi_b(0)=(2+4b)a^2.
\]
The left sides under consideration equal $\int\Phi_b(u)\,d\nu$.
The positive series \eqref{eq:h-expansion}, together with the convex
quadratic in this definition, shows that $\Phi_b$ is convex on $[0,1]$
for every $b\ge0$. Its endpoint difference is
\begin{equation}\label{eq:Phi-endpoints}
 \Phi_b(1)-\Phi_b(0)
 =-a^2-2a-2\log(1-a)-b(4a^3-a^4).
\end{equation}
For $a=9/10$ and $b\ge11/8$ this is negative. This can be checked with
rational arithmetic and $\log10<5/2$; a sharper enclosure is recorded in
Appendix~\ref{app:certificate}. Convexity therefore gives
$\Phi_b(x)\le\Phi_b(0)$ throughout $[0,1]$, proving
\eqref{eq:scalar-large-b}.

For $b=1/3$, expansion gives
\[
 \Phi_{1/3}(x)=\frac{10a^2}{3}-\frac{2a^3}{3}x
       +\frac{5a^4}{6}x^2
       +2\sum_{k\ge5}\frac{a^kx^{k-2}}{k}.
\]
As $x\ge x^2$ on $[0,1]$ and
$5a^4/6-2a^3/3>0$ for our value of $a$,
\[
 \Phi_{1/3}(x)\le\frac{10a^2}{3}
 +\left(\frac{5a^4}{6}-\frac{2a^3}{3}\right)x^2
 +2\sum_{k\ge5}\frac{a^kx^{k-2}}{k}.
\]
Integrating against $\nu$ and applying
Corollary~\ref{cor:normalized-moments}, with $c_4=I_4=2/3$, proves
\eqref{eq:scalar-first}. To justify the infinite sum uniformly in $B$,
first truncate at a fixed index. The remaining contribution is at most
$2\sum_{k>L}a^k/k$, since $n\int u^k\le1$. This tends to zero independently
of $B$, after which the finite number of moment limits may be taken.
\end{proof}

\begin{proposition}[Uniform tail budget]\label{prop:tail-budget}
For the ratio-nine chain above, as its smallest size tends to infinity,
\begin{equation}\label{eq:tail-budget}
 n_i\norm{P_i-1}_2^2\le\mathcal E_{9/10}+o(1)
\end{equation}
uniformly in $i$, in the length of the chain, and in the underlying set.
\end{proposition}
\begin{proof}
Apply Lemma~\ref{lem:scalar} in \eqref{eq:coupled-tail}. All terms with
$r\ge2$ are bounded without error, and only $r=1$ uses a moment limit.
If there is no tail, the result is immediate. Otherwise the required
moment error is uniform once $n_{i+1}$ tends to infinity. The constant
part is bounded by the infinite sum
\[
 a^2\sum_{r\ge1}w_r^2(2+4b_r)
 =2a^2V_0\sum_{r\ge1}\frac{w_r^2}{v_r}
 =2a^2\frac5{12}\left(\frac49+\frac16\right)
 =\frac{55a^2}{108}.
\]
The first-level correction is $w_1^2\Delta_a=\Delta_a/9$. This is exactly
$\mathcal E_a$ from \eqref{eq:E-intro}. Adding unused nonnegative terms
causes no difficulty for a finite tail.
\end{proof}

\subsection{Completion of the quantitative theorem}

\begin{proof}[Proof of Theorem~\ref{thm:constant}]
For large $N$, put $n_1=\lceil\log N\rceil$ and
$n_i=9^{i-1}n_1$ for $1\le i\le J$, where $J$ is maximal with
$n_J\le N$. Thus
\begin{equation}\label{eq:chain-length}
 J=\frac{\log N-\log\log N}{\log9}+O(1).
\end{equation}
Let $A_i$ be the first $n_i$ elements of $A$, and use the bounded test
$G$ in \eqref{eq:test-function}. The interval energy bound
\eqref{eq:energy-max} gives
\[
 \frac{E(A_i)+n_i^2}{2n_i^2}
 \le\frac{n_i}{3}+\frac12+\frac1{6n_i}.
\]
Together with Lemma~\ref{lem:prefix-pairing} and
Proposition~\ref{prop:tail-budget}, this implies, uniformly in $i$,
\[
 \Rea\ip{F_A}{(F_{A_i}/n_i)P_i}
 \ge1-\sqrt{\mathcal E_{9/10}/3}-o(1).
\]
Because $\norm G_\infty\le1$, summing and using
\eqref{eq:chain-length} yields
\[
 \La(A)\ge\Rea\ip{F_A}G
 \ge\left[\frac{9}{10\log9}
       \left(1-\sqrt{\mathcal E_{9/10}/3}\right)-o(1)\right]\log N.
\]
Every error used here depends only on $N$, not on the diameter or
arithmetic arrangement of $A$. Finally, Appendix~\ref{app:certificate}
proves with rational bounds that
\[
 \mathcal E_{9/10}<0.479084457020255861,
 \qquad c_*>0.2459209213577253.
\]
In particular $c_*>0.2459209$, as stated.
\end{proof}

\begin{remark}
The last theorem is an asymptotic leading-coefficient statement. It does
not assert the same decimal times $\log N$ with no error for every small
$N$, nor does it optimize the entire class of outer recursions. Its
frequency support still has a one-sided endpoint restriction. The
classical Landau bound in Appendix~\ref{app:landau} explains why
optimizing only within that final spectral class cannot reach $4/\pi^2$.
\end{remark}

\Needspace{7\baselineskip}
\section{Hereditary energy and repeated structural extraction}\label{sec:hered}

The structural proof needs an invariant that survives deletion. A
bound for the Fourier norm of an original set need not survive when a
subset is removed. An energy inequality formulated for \emph{every}
remaining subset does survive, and this is the purpose of the next
lemma.

\begin{lemma}[Hereditary energy gives an almost-covering]\label{lem:peeling}
Let $X$ be a nonempty finite subset of an abelian group, $|X|=R$, and
suppose that, for some $0<\eta\le1$,
\begin{equation}\label{eq:hereditary-energy}
 E(Y)\ge\eta\frac{|Y|^4}{R}
 \qquad(\varnothing\ne Y\subseteq X).
\end{equation}
For every $0<\delta<1$, there are pairwise disjoint subsets
$X_1,\ldots,X_J\subseteq X$ such that
\begin{align}
 \left|X\setminus\bigcup_{j=1}^JX_j\right|&<\delta R,
 &J&\le C\eta^{-C}\delta^{-C},\label{eq:peeling-cover}\\
 |X_j|&\ge c\eta^C\delta^C R,
 &|X_j+X_j|&\le C\eta^{-C}\delta^{-C}|X_j|.
 \label{eq:peeling-size}
\end{align}
All constants here are absolute.
\end{lemma}
\begin{proof}
Take the polynomial constants from Theorem~\ref{thm:BSG}. If the current
remainder $Y$ has $|Y|\ge\delta R$, then
\eqref{eq:hereditary-energy} gives
$E(Y)\ge\eta\delta|Y|^3$. BSG produces $Z\subseteq Y$ of size at least
$c(\eta\delta)^C|Y|$, and with doubling at most
$C(\eta\delta)^{-C}$. In particular the number removed is at least
$c(\eta\delta)^C\delta R$. Remove this $Z$ and repeat. The selected sets
are disjoint, every step removes a fixed parameter-dependent fraction
of the \emph{original} size $R$, and the process stops after a bounded
number of steps. Increasing the absolute exponent absorbs the additional
factor of $\delta$ in the size and step bounds. The only hypothesis
used at a later step is \eqref{eq:hereditary-energy} for its remainder.
\end{proof}

\begin{lemma}[A constant spectral budget gives hereditary energy]
\label{lem:spectral-hereditary}
Let $X$ be a nonempty finite subset of an abelian group, let $R=|X|$,
and put $S=\norm{F_X}_1$ on its compact dual. For every $Y\subseteq X$,
\begin{equation}\label{eq:spectral-hereditary}
 E(Y)\ge\frac{|Y|^4}{S^2R}.
\end{equation}
\end{lemma}
\begin{proof}
The claim is immediate for $Y=\varnothing$. Otherwise
$\ip{F_X}{F_Y}=|Y|$. Weighted Cauchy--Schwarz and then ordinary
Cauchy--Schwarz imply
\begin{align*}
 |Y|^2
 &\le\left(\int|F_X||F_Y|\right)^2\\
 &\le\left(\int|F_X|\right)
       \left(\int|F_X||F_Y|^2\right)\\
 &\le S\left(\int|F_X|^2\right)^{1/2}
          \left(\int|F_Y|^4\right)^{1/2}
 =S\sqrt R\sqrt{E(Y)}.
\end{align*}
Squaring and rearranging proves the result. In particular, this argument
uses no order on the group and applies on a finite cyclic group.
\end{proof}

Applied directly with $S=K\log N$, the last lemma still loses powers of
$\log N$. Its role below is to act \emph{after} a long integer direction
has been removed, when the residual spectral budget is a constant.
For separated blocks, we instead obtain \eqref{eq:hereditary-energy}
from bounded nonlinear amplification.

\Needspace{7\baselineskip}
\section{Bounded amplification and the energy of integer block labels}
\label{sec:majority}

\subsection{A positive density of balanced correlations}

\begin{lemma}[Dense mixed correlations]\label{lem:majority}
Let $F\in L^1(\T)$, $L=\norm F_1>0$, and let
$\Phi_1,\ldots,\Phi_r$ satisfy
\[
 \norm{\Phi_i}_\infty\le1,\qquad
 \Rea\ip F{\Phi_i}\ge\kappa>0.
\]
Put $T=L/\kappa\ge1$, and let $m$ be the smallest odd integer at least
$64T^2$. If $r\ge m$, then there is an odd integer
$k=2s-1$, $3\le k\le m$, such that at least $e^{-CT^2}r^k$ ordered
$k$-tuples of distinct indices satisfy
\begin{equation}\label{eq:good-correlations}
 \left|\ip F{\Phi_{i_1}\cdots\Phi_{i_s}
       \overline{\Phi_{i_{s+1}}\cdots\Phi_{i_k}}}\right|
       \ge e^{-CT^2}\kappa.
\end{equation}
No arithmetic assumption on the test functions is made.
\end{lemma}
\begin{proof}
For odd $m$, let $\Maj_m:\{-1,1\}^m\to\{-1,1\}$ be the sign of the
sum of its coordinates, and consider its multi-affine extension
\[
 P_m(t_1,\ldots,t_m)=
 \sum_{S\subseteq[m]}\widehat{\Maj_m}(S)\prod_{j\in S}t_j.
\]
This extension is bounded in absolute value by one throughout
$[-1,1]^m$: it is the expectation of Boolean majority under independent
signs whose expectations are the prescribed $t_j$. Only odd degrees
occur. Each first-degree coefficient equals
\[
 a_m=2^{-(m-1)}\binom{m-1}{(m-1)/2}\ge(2m)^{-1/2}.
\]
Indeed the other $m-1$ signs contribute to the correlation with one
coordinate exactly when they sum to zero. The displayed lower bound
follows from the elementary central-binomial estimate. Boolean Parseval
and Cauchy--Schwarz also give
\begin{equation}\label{eq:walsh-mass}
 B_m:=\sum_{S\subseteq[m]}|\widehat{\Maj_m}(S)|\le2^{m/2}.
\end{equation}

To handle complex functions without losing the balance of conjugates,
work on $\T^2$ and put
\[
 H(x,t)=2\Rea(e(t)F(x)),\qquad
 \psi_i(x,t)=\Rea(e(t)\Phi_i(x)).
\]
Then $|\psi_i|\le1$ and direct integration in $t$ gives
\begin{equation}\label{eq:phase-lift}
 \norm H_1=\frac4\pi L=:L_*,\qquad
 \ip H{\psi_i}=\Rea\ip F{\Phi_i}\ge\kappa.
\end{equation}
For any injection $\iota:[m]\to[r]$, pair $H$ with
$P_m(\psi_{\iota(1)},\ldots,\psi_{\iota(m)})$. The full pairing has
absolute value at most $L_*$. Its linear contribution is at least
$A_m\kappa$, where
\[
 A_m=ma_m\ge\sqrt{m/2}\ge2L_*/\kappa.
\]
Average over all injections. For an odd degree $k$, let
\[
 J_k=\E_{i_1,\ldots,i_k\,\mathrm{distinct}}
           \ip H{\psi_{i_1}\cdots\psi_{i_k}}.
\]
Every monomial of degree $k$ has the same averaged pairing $J_k$.
The higher terms must offset at least $A_m\kappa-L_*\ge A_m\kappa/2$.
By \eqref{eq:walsh-mass}, some odd $k\ge3$ therefore satisfies
\begin{equation}\label{eq:Jk}
 |J_k|\ge\frac{A_m}{2B_m}\kappa=:\tau\kappa,
 \qquad \tau\ge e^{-Cm}.
\end{equation}

Write $k=2s-1$. Expanding the real parts and integrating in the extra
phase yields the exact identity
\begin{equation}\label{eq:balanced-identity}
 \ip H{\prod_{j=1}^k\psi_{i_j}}
 =2^{1-k}\Rea\sum_{\substack{U\subseteq[k]\\|U|=s}}
  \ip F{\prod_{j\in U}\Phi_{i_j}
                   \prod_{j\notin U}\overline{\Phi_{i_j}}}.
\end{equation}
Only products with one more unconjugated factor survive. After
averaging, all the terms on the right agree by permutation symmetry.
Since $2^{1-k}\binom{k}{s}\le1$, \eqref{eq:Jk} implies
\[
 \E_{\mathrm{distinct}}
 \left|\ip F{\Phi_{i_1}\cdots\Phi_{i_s}
       \overline{\Phi_{i_{s+1}}\cdots\Phi_{i_k}}}\right|
 \ge\tau\kappa.
\]
Each absolute pairing is at most $L=T\kappa$. Hence at least a
$\tau/(2T)$ fraction of the distinct tuples have absolute pairing at
least $\tau\kappa/2$. Finally,
\[
 \frac{(r)_k}{r^k}\ge\frac{k!}{k^k}\ge e^{-k}
 \qquad(r\ge k),
\]
where $(r)_k=r(r-1)\cdots(r-k+1)$. Since $k\le m\ll T^2$, all fixed
factors and the factor $T$ can be absorbed into $e^{CT^2}$. This proves
\eqref{eq:good-correlations} with an absolute $C$.
\end{proof}

The distinction between a single nonzero mixed correlation and a positive
density of such correlations is essential. The latter can be counted
as additive relations. The argument has also retained all higher terms
of the bounded polynomial; it has not presumed that a cubic correction
alone captures the obstruction.

\subsection{Finite-band certificates for actual integer blocks}

Let $C$ have the exact representation
\eqref{eq:terminal-packet}, initially without the density conditions
in \eqref{eq:terminal-packet-conditions}. We require only $D\ge1$,
$C_i\ne\varnothing$, and separation $|b_i-b_j|>4D$.

\begin{lemma}[A finite-band certificate]\label{lem:band}
For every $B\subseteq[-D,D]\cap\Z$ there is a trigonometric polynomial
$\psi_B$ such that
\[
 \norm{\psi_B}_\infty\le1,\qquad
 \supp\widehat{\psi_B}\subseteq[-2D,2D],\qquad
 \ip{F_B}{\psi_B}=\frac13\La(B).
\]
\end{lemma}
\begin{proof}
Let $\mathcal F_D$ be the Fej\'er kernel whose coefficients are
$(1-|n|/D)_+$, and put $V_D=2\mathcal F_{2D}-\mathcal F_D$.
Its coefficients equal one on $[-D,D]$, vanish outside $(-2D,2D)$,
and $\norm{V_D}_1\le3$. Define the phase of $F_B$ to be zero at its
zeros, and put
\[
 \psi_B=\frac13 V_D*\frac{F_B}{|F_B|}.
\]
Convolution gives the $L^\infty$ and support assertions. Since $V_D$ is
real and even and has coefficient one at every frequency of $B$,
\[
 \ip{F_B}{\psi_B}
 =\frac13\ip{V_D*F_B}{F_B/|F_B|}
 =\frac13\int|F_B|.
\]
The empty-set case is covered by the zero function.
\end{proof}

Take $\Phi_i=e(b_it)\psi_{C_i}$. Its support lies in
$b_i+[-2D,2D]$. The separation of the centres implies that no other
block $b_j+C_j$ meets this support. Consequently
\begin{equation}\label{eq:ambient-certificate}
 \ip{F_C}{\Phi_i}=\frac13\La(C_i).
\end{equation}
Let
\begin{equation}\label{eq:label-def}
 \ell_i=\lfloor b_i/D\rfloor,
 \quad \Lambda=\{\ell_i:1\le i\le R\},
 \quad\kappa=\frac13\min_i\La(C_i),
 \quad T=\La(C)/\kappa.
\end{equation}
The labels are distinct and $T\ge1$ by \eqref{eq:ambient-certificate}.

\begin{proposition}[Hereditary energy of the labels]\label{prop:packet-hered}
For every nonempty $S\subseteq\Lambda$,
\begin{equation}\label{eq:packet-hered}
 E(S)\ge e^{-CT^2}\frac{|S|^4}{R}.
\end{equation}
The parameter $T$ refers to the full original target $C$ and is the same
for every $S$.
\end{proposition}
\begin{proof}
Write $r=|S|$, and let $m$ be as in Lemma~\ref{lem:majority}. If $r<m$,
then $E(S)\ge r^2$ suffices: because $R\ge r$ and $m\ll T^2$, an
absolute enlargement of $C$ ensures $e^{-CT^2}r^2\le R$.
Assume therefore that $r\ge m$. Apply Lemma~\ref{lem:majority} to the
certificates indexed by $S$, keeping $F_C$ as target. For some
$k=2s-1\ll T^2$, at least $e^{-CT^2}r^{2s-1}$ ordered distinct tuples
have nonzero balanced correlation.

The Fourier support of such a balanced product lies in
\[
 b_{i_1}+\cdots+b_{i_s}-b_{i_{s+1}}-\cdots-b_{i_k}
                    +[-2kD,2kD].
\]
Nonzero correlation with $F_C$ forces this interval to contain a point
of one of the original blocks. Hence, for some $j\in\{1,\ldots,R\}$,
\[
 |b_{i_1}+\cdots+b_{i_s}-b_{i_{s+1}}-\cdots-b_{i_k}-b_j|
 \le(2k+1)D.
\]
The witness $j$ is not required to have $\ell_j\in S$. Writing
$b_i=D\ell_i+r_i$, $0\le r_i<D$, gives an integer $t$ with
\begin{equation}\label{eq:carry-relation}
 \ell_{i_1}+\cdots+\ell_{i_s}
 -\ell_{i_{s+1}}-\cdots-\ell_{i_k}-\ell_j=t,
 \qquad |t|\le4k.
\end{equation}
Indeed the rounding error contains $k+1$ remainders, so $(2k+1)+(k+1)$
provides a sufficient bound for $k\ge3$. Choose one witness for each
good tuple. Different tuples still give different ordered solutions
when the witness is included. With $r_s=\one_S^{*s}$, their count yields
\begin{equation}\label{eq:ambient-relations}
 \sum_{|t|\le4k}\sum_x r_s(x)
           (\one_S^{*(s-1)}*\one_\Lambda)(x+t)
 \ge e^{-CT^2}r^{2s-1}.
\end{equation}
Changing $t$ to $-t$ if necessary gives exactly this convention.

By Fourier inversion and then H\"older, each inner shifted pairing is
bounded above by
\begin{align*}
 \int_\T|F_S|^{2s-1}|F_\Lambda|
 &\le r^{2s-4}\int_\T|F_S|^3|F_\Lambda|\\
 &\le r^{2s-4}E(S)^{3/4}E(\Lambda)^{1/4}\\
 &\le r^{2s-4}E(S)^{3/4}R^{3/4}.
\end{align*}
Here $s\ge2$, and the last inequality uses the elementary
$E(\Lambda)\le R^3$. Combining with \eqref{eq:ambient-relations},
absorbing $8k+1$ into the exponential constant, and taking the power
$4/3$ proves \eqref{eq:packet-hered}.
\end{proof}

In particular, the high-order correlations have been reduced to ordinary
fourth-order energy. No assertion that the floor map is a Freiman
isomorphism is used: the bounded integer error in
\eqref{eq:carry-relation} is counted rather than discarded.

\begin{theorem}[Almost-covering for dense separated blocks]
\label{thm:packet-cover}
Fix $K_0\ge1$, $0<\beta\le1$, $\Delta\ge1$, and $0<\eps<1$.
If a set $C$ satisfies \eqref{eq:terminal-packet}--
\eqref{eq:terminal-packet-conditions}, $|C|=n\ge2$, and
$\La(C)\le K_0\log n$, then all but at most $\eps n$ points of $C$
are covered by boundedly many pairwise disjoint subsets, each of size
at least $cn$ and doubling at most $C_1$. The positive constants and
the number of subsets depend only on $K_0,\beta,\Delta,\eps$.
\end{theorem}
\begin{proof}
The weighted MPS inequality gives
\[
 \kappa\ge c\log(M+1)\ge c\beta\log n,
 \qquad T\le CK_0/\beta.
\]
Apply Lemma~\ref{lem:peeling} to \eqref{eq:packet-hered}, with
$\delta=\eps/(6\Delta)$. The omitted labels number less than
$\delta R$. Since $|C_i|\le2D+1\le3\Delta M$ and $n\ge MR$,
their original blocks contain at most
\[
 (2D+1)\delta R\le3\Delta\delta n\le\eps n/2
\]
points. For each selected label set $S_j$, take $B_j$ to be the union
of all its original blocks. The $B_j$ are disjoint and
$|B_j|\ge M|S_j|$. Also $n\le3\Delta MR$, so the lower size bound for
$S_j$ gives $|B_j|\ge cn$.

If $|S_j+S_j|\le Q|S_j|$, each point in a corresponding integer block
has the form $D\ell+u$ with $u\in[-D,2D]\cap\Z$. Thus
\[
 |B_j+B_j|\le(6D+1)|S_j+S_j|
 \le7\Delta Q|B_j|.
\]
The label constants are bounded in the stated parameters, and so are
all the constants in the conclusion. In particular no estimate for
$\La$ of a deleted-block remainder was required.
\end{proof}

\Needspace{7\baselineskip}
\section{Integer boundaries and long arithmetic blocks}\label{sec:boundary}

The second local argument converts a logarithmic budget into a constant
spectral budget. The relevant object is an integer-valued boundary on a
cyclic group times $\Z$, rather than a collection of independently
chosen analytic vectors.

\subsection{Extracting the boundary norm}

Let $q,M\ge1$ and let
$\varnothing\ne\mathcal C\subset(\Z/q\Z)\times\Z$ be finite, of size
$V$. Initially no aspect condition is imposed. Residues $b$ are always
represented by $0\le b<q$. Put $A=A_M(\mathcal C)$ as in
\eqref{eq:terminal-inflation}, and define
\begin{align}
 h(b,u)&=\one_{\mathcal C}(b,u)-\one_{\mathcal C}(b,u-1),
 &H&=\sum_{b,u}|h(b,u)|^2\le2V,\label{eq:boundary-def}\\
 \mathscr H(j,s)&=\sum_{b,u}h(b,u)e(bj/q+us),
 &S&=\frac1q\sum_{j=0}^{q-1}\int_\T|\mathscr H(j,s)|\,ds.
 \label{eq:boundary-norm}
\end{align}
The function $h$ is nonzero, takes values in $\{-1,0,1\}$, and has zero
sum in $u$ for every $b$. Fourier inversion and Parseval give
\begin{equation}\label{eq:S-bounds}
 1\le S\le\sqrt H.
\end{equation}

\begin{proposition}[Uniform boundary extraction]\label{prop:boundary}
For every $q,M,\mathcal C$ as above,
\begin{equation}\label{eq:boundary-extraction}
 \La(A_M(\mathcal C))
 \ge\frac S\pi\log\frac{MS}{\sqrt H}-4S.
\end{equation}
For fixed $q,\mathcal C$, as $M\to\infty$ one also has
\begin{equation}\label{eq:boundary-asymptotic}
 \La(A_M(\mathcal C))=\frac S\pi\log M+O_{q,\mathcal C}(1).
\end{equation}
\end{proposition}
\begin{proof}
The finite geometric series for each cell gives the exact identity
\begin{equation}\label{eq:boundary-identity}
 (1-e(q\theta))F_A(\theta)
 =\sum_{b,u}h(b,u)e(b\theta+qMu\theta).
\end{equation}
Partition the circle by writing $\theta=(j+t)/q$, with
$0\le j<q$ and $-1/2\le t<1/2$. Define
\[
 W(t,s)=\frac1q\sum_{j=0}^{q-1}
       \left|\sum_{b,u}h(b,u)e(b(j+t)/q+us)\right|.
\]
Equation \eqref{eq:boundary-identity} yields
\begin{equation}\label{eq:boundary-integral}
 \La(A)=\int_{-1/2}^{1/2}
            \frac{W(t,Mt)}{2|\sin\pi t|}\,dt.
\end{equation}
The values at the isolated zeros of the denominator are interpreted
by the original identity and do not affect the integral.

Let
\[
 W_0(s)=W(0,s),\qquad
 Q(s)=\left(\sum_b\left|\sum_uh(b,u)e(us)\right|^2\right)^{1/2}.
\]
For fixed $s$, the difference inside $W(t,s)$ and $W_0(s)$ is a
Fourier sum in $j$ with coefficients
$(e(bt/q)-1)\sum_uh(b,u)e(us)$. The triangle inequality, discrete
Parseval, and $|e(bt/q)-1|\le2\pi|t|$ prove
\begin{equation}\label{eq:boundary-freezing}
 |W(t,s)-W_0(s)|\le2\pi|t|Q(s).
\end{equation}
Since $h$ is real, $W_0,Q$ are even and one-periodic, and
$W(-t,-s)=W(t,s)$ after conjugation and reindexing $j$.
Furthermore
\begin{equation}\label{eq:boundary-averages}
 \int_0^1W_0(s)\,ds=S,\qquad
 \int_0^1Q(s)\,ds\le\sqrt H.
\end{equation}
The second assertion follows from Cauchy--Schwarz and Parseval in $s$.

Restrict \eqref{eq:boundary-integral} to $|t|\le a$, where
$0<a\le1/2$. Since $2t\le\sin\pi t\le\pi t$ for $0\le t\le1/2$,
\eqref{eq:boundary-freezing} gives
\begin{equation}\label{eq:boundary-small-arc}
 \La(A)\ge\frac1\pi\int_0^{Ma}\frac{W_0(u)}u\,du
              -\pi(a+M^{-1})\sqrt H.
\end{equation}
For the error term, split $[0,Ma]$ into complete unit periods and one
remaining interval; nonnegativity and periodicity bound its integral
of $Q$ by $(Ma+1)\int_0^1Q$.

A nonnegative one-periodic function of mean $S$ satisfies
\begin{equation}\label{eq:periodic-harmonic}
 \int_0^X\frac{W_0(u)}u\,du\ge S(\log X-2)\qquad(X\ge1).
\end{equation}
To see this, discard $[0,1]$ and the final incomplete interval. Each
complete interval $[j,j+1]$, $j\ge1$, contributes at least $S/(j+1)$.
The elementary harmonic-sum bound proves \eqref{eq:periodic-harmonic}.
Now choose $a=S/(4\sqrt H)$, which lies in $(0,1/4]$ by
\eqref{eq:S-bounds}. If $MS/\sqrt H\ge4$, then $Ma\ge1$ and the error
in \eqref{eq:boundary-small-arc} is at most $\pi S/2$. Hence
\[
 \La(A)\ge\frac S\pi\log\frac{MS}{\sqrt H}
       -\left(\frac{\log4+2}{\pi}+\frac\pi2\right)S,
\]
which is stronger than \eqref{eq:boundary-extraction}. If
$MS/\sqrt H<4$, the right side of \eqref{eq:boundary-extraction} is
negative, so that inequality is automatic.

For \eqref{eq:boundary-asymptotic}, hold $q,\mathcal C$ fixed. On the
whole interval $[-1/2,1/2]$, the cruder bound
$|W(t,s)-W_0(s)|\le2\pi|t|\sum|h|$ has an integrable quotient by
$|\sin\pi t|$. Replacing $1/(2|\sin\pi t|)$ by $1/(2\pi|t|)$ also
has bounded integrated error. Thus
\[
 \La(A)=\frac1\pi\int_0^{M/2}\frac{W_0(u)}u\,du+O_{q,\mathcal C}(1).
\]
The integral near zero is finite: $W_0(0)=0$, because each vertical
sum of $h$ is zero, and $W_0$ is Lipschitz. The periodic function
$W_0-S$ has a bounded primitive; integration by parts on $[1,M/2]$
therefore shows that its contribution divided by $u$ is bounded.
This proves \eqref{eq:boundary-asymptotic}.
\end{proof}

\begin{corollary}[A constant boundary budget]\label{cor:boundary-budget}
Suppose $M\ge V$, $n=MV$ is sufficiently large absolutely, and
$\La(A_M(\mathcal C))\le K_0\log n$, with $K_0\ge1$. Then
$S\le CK_0$. More generally, $M\ge V^{1/2+\sigma}$ with fixed
$0<\sigma\le1/2$ gives $S\le CK_0/\sigma$ once $n$ is sufficiently
large in terms of $\sigma$.
\end{corollary}
\begin{proof}
Since $H\le2V$ and $S\ge1$, the first aspect condition gives
\[
 \log\frac{MS}{\sqrt H}\ge\frac12\log M-\frac12\log2,
 \qquad\log n\le2\log M.
\]
For large $M$, Proposition~\ref{prop:boundary} implies
$\La(A_M(\mathcal C))\ge(S/(4\pi))\log M$, proving the first assertion.
The condition $M\ge V$ ensures that $M\to\infty$ as $n\to\infty$.
For the second assertion use
\[
 \log(M/\sqrt V)\ge\frac{2\sigma}{1+2\sigma}\log M,
 \qquad\log n\le3\log M,
\]
and absorb the fixed negative multiple of $S$ once more.
\end{proof}

\subsection{Integer-valuedness creates shared boundary heights}

\begin{lemma}[Common heights and residual spectra]\label{lem:heights}
Let $u_1<\cdots<u_v$ be the heights at which $h(\cdot,u)$ is nonzero.
Then
\begin{equation}\label{eq:height-count}
 v\le e^{CS}.
\end{equation}
Between consecutive heights, $\mathcal C$ is a slab
$X_j\times[u_j,u_{j+1})$, with possibly empty
$X_j\subseteq\Z/q\Z$. The nonempty slabs partition $\mathcal C$, and
\begin{equation}\label{eq:residue-spectrum}
 \norm{F_{X_j}}_{L^1(\widehat{\Z/q\Z})}\le vS.
\end{equation}
\end{lemma}
\begin{proof}
For fixed $a\in\{0,\ldots,q-1\}$, apply
Theorem~\ref{thm:MPS} in the vertical variable to
\[
 \mathscr H(a,s)=\sum_{j=1}^v
 \left(\sum_bh(b,u_j)e(ba/q)\right)e(u_js).
\]
Some displayed coefficients may vanish; the weighted inequality remains
valid with the original indices. Averaging over $a$ gives
\[
 S\ge\cMPS\sum_{j=1}^v\frac1j
       \norm{F_{h(\cdot,u_j)}}_{L^1(\widehat{\Z/q\Z})}.
\]
Each underlying coefficient function is nonzero and integer-valued, so
its Fourier algebra norm is at least one by inversion. Hence
$S\ge\cMPS\sum_{j\le v}1/j$, proving \eqref{eq:height-count}.

Extraction of the coefficient at height $u_j$ is contractive:
\[
 \frac1q\sum_a\left|\int_\T\mathscr H(a,s)e(-u_js)\,ds\right|
 \le S.
\]
Thus each boundary-height function has finite-group norm at most $S$.
Telescoping the definition of $h$ shows that, on the vertical interval
$[u_j,u_{j+1})$, the indicator of $\mathcal C$ is the residue indicator
\[
 \one_{X_j}(b)=\sum_{i\le j}h(b,u_i).
\]
It is zero before the first and after the last height, because
$\mathcal C$ is finite. This proves the slab partition and, by the
triangle inequality, \eqref{eq:residue-spectrum}.
\end{proof}

It is important that the values of $h$ are integers. Arbitrarily small
nonzero real coefficients would not impose a cost of at least one at a
new height. It is also important that the residue group remains cyclic:
replacing it by independent free coordinates would erase carries.

\begin{lemma}[Lifting doubling with carries]\label{lem:carry-lift}
Let $X\subseteq\{0,\ldots,q-1\}$ and
$P=X+q\{a,a+1,\ldots,a+L-1\}$, where $L\ge1$. If
$|X+X|_{\Z/q\Z}\le Q|X|$, then
\begin{equation}\label{eq:carry-doubling}
 |P+P|\le2Q|P|.
\end{equation}
\end{lemma}
\begin{proof}
For a residue sum write $x_1+x_2=r+q\kappa$ with
$0\le r<q$ and $\kappa\in\{0,1\}$. The sum of two vertical coordinates
in $[a,a+L)$ ranges from $2a$ to $2a+2L-2$. Allowing either carry
produces at most $2L$ values for that residue. Therefore
$|P+P|\le2L|X+X|_{\Z/q\Z}\le2QL|X|=2Q|P|$.
\end{proof}

\begin{theorem}[Almost-covering for long block inflations]
\label{thm:block-cover}
Fix $K_0\ge1$ and $0<\eps<1$. Let $C=A_M(\mathcal C)$ with
$M\ge V$, $n=MV$ sufficiently large, and $\La(C)\le K_0\log n$.
Then all but at most $\eps n$ points of $C$ are covered by boundedly
many disjoint subsets, each of size at least $cn$ and doubling at most
$C_1$. All constants depend only on $K_0,\eps$.
\end{theorem}
\begin{proof}
Corollary~\ref{cor:boundary-budget} and Lemma~\ref{lem:heights} give
$S\le CK_0$, $v\le e^{CK_0}$, and a partition into at most $v$
nonempty arithmetic slabs
\[
 P_j=X_j+q\{a_j,\ldots,a_j+L_j-1\}.
\]
Here $L_j=M(u_{j+1}-u_j)$, and the finite-group spectral norm of each
$X_j$ is at most $vS$, hence bounded only in terms of $K_0$.

Discard slabs of size less than $\eps n/(4v)$, losing at most
$\eps n/4$ points. For a remaining residue set $X_j$, apply
Lemma~\ref{lem:spectral-hereditary} and then
Lemma~\ref{lem:peeling}, with $\delta=\eps/4$. This covers all but a
$\delta$ proportion of its residues by boundedly many disjoint subsets
$X_{j,l}$, each of positive parameter-dependent relative size and
bounded doubling in $\Z/q\Z$. Lift each one over the entire vertical
interval of the slab. The lifted pieces are disjoint,
Lemma~\ref{lem:carry-lift} bounds their integer doubling, and the
omitted residue proportion equals the omitted point proportion.
Since every retained slab has size at least $\eps n/(4v)$, each
lifted piece is linear in $n$. The additional omitted mass is at most
$\eps n/4$. This proves the stated conclusion, with room in the error
budget.
\end{proof}

\begin{remark}[A stronger representation inside this terminal model]
Green--Sanders~\cite{GreenSanders} gives an additional description here:
$\one_C$ is an exact signed sum of $O_{K_0}(1)$ arithmetic-progression
indicators. Indeed each $\one_{X_j}$ has bounded finite-group algebra
norm and is a bounded signed sum of coset indicators. A coset in
$\Z/q\Z$ has representatives
$r+d\{0,\ldots,q/d-1\}$, $d\mid q$, $0\le r<d$. Its lift over a
vertical interval $[a,a+L)$ is the single progression
\[
 r+qa+d\{0,\ldots,(q/d)L-1\}.
\]
Lifting the signed identities slab by slab gives the assertion. This
observation is not needed for Theorem~\ref{thm:inverse}. Nor does it
bound the length of each progression by $|C|$: large progressions in
a signed representation may cancel. Thus it should not be identified
with the full signed-progression conjecture in \cite{BG}.
\end{remark}

\Needspace{7\baselineskip}
\section{A depth-independent Sidon budget and the inverse theorem}
\label{sec:gluing}

The local covering theorems are not sufficient by themselves for an
arbitrarily deep assembly. A multiplicative loss at every split could
vanish along a long path. We use a weaker coefficient exponent with
constant exactly one, and then sum over all leaves at once.

\subsection{A coefficient inequality with constant one}

\begin{lemma}[Constant-one Sidon aggregation]\label{lem:sidon20}
Let $R$ be a finite additive Sidon set in the dual of a compact abelian
group. For arbitrary complex coefficients,
\begin{equation}\label{eq:sidon20}
 \left\|\sum_{r\in R}a_r\chi_r\right\|_1
 \ge\left(\sum_{r\in R}|a_r|^{20}\right)^{1/20}.
\end{equation}
The Haar measure is normalized. The exponent $20$ is not asserted to
be optimal.
\end{lemma}
\begin{proof}
The zero polynomial is immediate. Divide by a coefficient of maximum
absolute value and multiply by the conjugate of its character. This
preserves both sides up to the same factor, preserves the additive
Sidon property, and leaves a polynomial $1+g$ whose other coefficients
have absolute value at most one. Let their squared absolute values
sum to $t$. The case $t=0$ is immediate, so assume $t>0$.

Orthogonality and the Sidon condition give
\begin{equation}\label{eq:sidon-moments}
 \E g=\E g^2=0,\qquad
 \E|g|^2=t,\qquad \E|g|^4\le2t^2.
\end{equation}
For example, $\E g^2=0$ follows because the sum of two nonzero shifted
labels cannot equal the sum $0+0$ of the distinguished label. For the
fourth moment only the two trivial pairings survive, with the all-equal
terms counted only once. In particular, if $y=\Ima g$, then
$\E y^2=t/2$.

For every $z\in\C$ one has
\begin{equation}\label{eq:pointwise-excess}
 |1+z|-1-\Rea z\ge\frac{(\Ima z)^2}{2(1+|z|)}.
\end{equation}
When $\Ima z\ne0$, rationalize the left side and bound its denominator
$|1+z|+1+\Rea z$ above by $2(1+|z|)$; when $\Ima z=0$, the right
side vanishes and the inequality is automatic. Weighted
Cauchy--Schwarz and \eqref{eq:sidon-moments} yield
\begin{align}
 \E|1+g|
 &\ge1+\frac{(\E y^2)^2}{2\E[y^2(1+|g|)]}\notag\\
 &\ge1+\frac{t}{4+8\sqrt{2t}}.
 \label{eq:sidon-small-t}
\end{align}
For the second line, use
$\E y^2|g|\le\E|g|^3\le\sqrt{\E|g|^2\E|g|^4}
\le\sqrt2\,t^{3/2}$.

If $t\le2$, \eqref{eq:sidon-small-t} is at least
$1+t/20\ge(1+t)^{1/20}$. If $t\ge2$, apply $L^1$--$L^2$--$L^4$
interpolation to the whole Sidon polynomial:
\[
 \E|1+g|\ge
 \frac{\norm{1+g}_2^3}{\norm{1+g}_4^2}
 \ge\sqrt{\frac{1+t}{2}}
 \ge(1+t)^{1/20}.
\]
The last inequality holds for $1+t\ge3$. Finally, since each
normalized coefficient has modulus at most one, its twentieth power
is no larger than its square. Thus the normalized right side of
\eqref{eq:sidon20} is at most $(1+t)^{1/20}$, proving the lemma.
\end{proof}

\begin{proposition}[Fibre and leaf budgets]\label{prop:leaf-budget}
At any internal vertex of Definition~\ref{def:class},
\begin{equation}\label{eq:fibre-budget}
 \La(A_v)^{20}\ge\sum_{r\in R_v}\La(A_{v,r})^{20}.
\end{equation}
Consequently, for an assembly tree with leaves $A_l$,
\begin{equation}\label{eq:leaf-budget}
 \sum_l\La(A_l)^{20}\le\La(A)^{20}.
\end{equation}
If $|A|=N$ and $\La(A)\le K\log N$, then for every fixed
$0<\theta<1$,
\begin{align}
 \#\{l\}&\le(K\log N)^{20},\label{eq:leaf-count}\\
 \#\{l:|A_l|\ge N^\theta\}
 &\le\left(\frac{K}{\cMPS\theta}\right)^{20},
 \label{eq:large-leaf-count}\\
 \sum_{|A_l|<N^\theta}|A_l|
 &\le N^\theta(K\log N)^{20}=o_K(N).
 \label{eq:small-leaf-mass}
\end{align}
\end{proposition}
\begin{proof}
Affine normalization preserves all relevant norms. For the normalized
vertex, write $\xi=(j+x)/q_v$. At fixed $x$ its Fourier polynomial is
\[
 F_{A_v}((j+x)/q_v)
 =\sum_{r\in R_v} e(rx/q_v)F_{A_{v,r}}(x)e(rj/q_v).
\]
Apply Lemma~\ref{lem:sidon20} on the finite cyclic group in $j$, and
then integrate over $x\in[0,1]$. Minkowski's inequality in the finite
space $\ell^{20}(R_v)$ gives
\begin{align*}
 \La(A_v)
 &\ge\int_0^1
       \left(\sum_r|F_{A_{v,r}}(x)|^{20}\right)^{1/20}\,dx\\
 &\ge\left(\sum_r\La(A_{v,r})^{20}\right)^{1/20}.
\end{align*}
This proves \eqref{eq:fibre-budget}. Induction on the finite tree proves
\eqref{eq:leaf-budget}, without a factor depending on its depth.

Every nonempty leaf has norm at least one, proving
\eqref{eq:leaf-count}. The weighted MPS bound gives
$\La(A_l)\ge\cMPS\log(|A_l|+1)\ge\cMPS\theta\log N$
for a large leaf. Substitute this in \eqref{eq:leaf-budget} to obtain
\eqref{eq:large-leaf-count}. Finally the number of small leaves times
$N^\theta$ bounds their total cardinality, proving
\eqref{eq:small-leaf-mass}. Their affine images in $A$ are disjoint, so
these cardinalities are also their actual masses in the root set.
\end{proof}

\subsection{Proof of the almost-covering theorem}

\begin{proof}[Proof of Theorem~\ref{thm:inverse} for $A\in\mathfrak G_{\beta,\Delta}$]
Let
\[
 B_0=\max\left\{1,\left(\frac{2K}{\cMPS}\right)^{20}\right\}.
\]
Use Proposition~\ref{prop:leaf-budget} with $\theta=1/2$. There are
at most $B_0$ leaves of size at least $\sqrt N$. All other leaves
together have size at most $\sqrt N(K\log N)^{20}$. Choose $N_0$ large
enough that this quantity is at most $\eps N/4$ for $N\ge N_0$, and
discard these small leaves.

Among the remaining leaves discard any of size less than
$\eps N/(4B_0)$. Since there are at most $B_0$ of them, this removes
at most another $\eps N/4$ points. Every retained leaf has
\begin{equation}\label{eq:retained-leaf}
 n_l:=|A_l|\ge\frac{\eps N}{4B_0},\qquad
 \La(A_l)\le\La(A)\le K\log N\le2K\log n_l
\end{equation}
for sufficiently large $N$, depending only on $K,\eps$.
Enlarge $N_0$ also so that these leaves exceed the size thresholds in
Theorem~\ref{thm:block-cover}. They are not singletons. Their affine
normalizations are one of the two terminal geometries.

Apply Theorem~\ref{thm:packet-cover} or
Theorem~\ref{thm:block-cover} to each retained leaf, with norm parameter
$2K$ and local error $\eps/2$. The sum of these local exceptional
masses is at most $\eps N/2$. There are at most $B_0$ leaves and a
bounded number of pieces per leaf. Each piece is a fixed
parameter-dependent fraction of its leaf, and by
\eqref{eq:retained-leaf} therefore a fixed fraction of $N$.

Lift the pieces through the affine maps along their paths to the root.
These maps preserve cardinality and sumset doubling. Distinct lifted
pieces remain disjoint, because the child images partition each parent.
The total omitted mass is at most
$\eps N/4+\eps N/4+\eps N/2=\eps N$. This proves
\eqref{eq:cover-conclusion}, with constants depending only on
$K,\beta,\Delta,\eps$. The Freiman-model extension will be proved in
Section~\ref{sec:transport}.
\end{proof}

\begin{proof}[Proof of Corollary~\ref{cor:main-hered}]
Apply Theorem~\ref{thm:inverse} with error $\eps=\delta/4$. Denote its
number of pieces by $J$ and its doubling bound by $Q$. If
$Y\subseteq A$ has $|Y|\ge\delta N$, then at least $\delta N/2$ of
its points lie in the pieces. Thus, for some $j$,
$X=Y\cap B_j$ has size at least $\delta N/(2J)$. Since
$|X+X|\le|B_j+B_j|\le QN$, \eqref{eq:energy-doubling} gives
\[
 E(Y)\ge E(X)\ge\frac{|X|^4}{QN}
       \ge\frac{\delta^4}{16QJ^4}N^3
       \ge\frac{\delta^4}{16QJ^4}|Y|^3.
\]
The lower bound for the constant depends only on the stated parameters.
\end{proof}

\begin{remark}[Why the hereditary formulation is natural]
For a uniform family of sets, the almost-covering conclusion for every
fixed $\eps>0$ is qualitatively equivalent to positive normalized
energy for every subset of fixed positive density. One direction is
the preceding proof. For the converse, while a remainder has at least
$\eps N$ elements, use its energy lower bound and BSG to remove at
least a fixed positive fraction of the original $N$. This terminates
in boundedly many steps. No monotonicity of the Fourier $L^1$ norm under
deletion is part of either argument.
\end{remark}

\Needspace{7\baselineskip}
\section{High-order modelling and small exceptional perturbations}
\label{sec:transport}

An order-two Freiman isomorphism preserves additive energy and sumset
cardinality, but does not generally preserve a Fourier $L^1$ norm.
For the transport used here we retain sufficiently many moments. This
also provides a bound for every corresponding subset simultaneously.

\begin{definition}[Freiman isomorphism of order $h$]
A bijection $\phi:A\to\widetilde A$ between finite subsets of abelian
groups is a Freiman isomorphism of order $h$ if
\[
 a_1+\cdots+a_h=b_1+\cdots+b_h
 \quad\Longleftrightarrow\quad
 \phi(a_1)+\cdots+\phi(a_h)=\phi(b_1)+\cdots+\phi(b_h)
\]
for all choices with repetitions allowed. For nonempty sets, padding
both sides with the same element shows that this condition also
preserves equal-length sum relations of every order at most $h$.
\end{definition}

\begin{lemma}[Uniform transport of the norm]\label{lem:transport}
Let $h\ge2$, and let $\phi:A\to\widetilde A$ be a Freiman isomorphism
of order $h$. For every $C\subseteq A$, with $m=|C|$,
\begin{equation}\label{eq:transport}
 |\La(C)-\La(\phi(C))|\le\frac{4m}{\pi(2h+1)}.
\end{equation}
Here the norms are taken on the corresponding compact dual groups.
Moreover,
\begin{equation}\label{eq:exact-transport}
 E(C)=E(\phi(C)),\qquad |C+C|=|\phi(C)+\phi(C)|.
\end{equation}
\end{lemma}
\begin{proof}
The empty-set case is immediate. The relation hypothesis and
orthogonality give
\[
 \int|F_C|^{2j}=\int|F_{\phi(C)}|^{2j}\qquad(1\le j\le h).
\]
For $-1\le t\le1$, the absolutely convergent Chebyshev expansion is
\begin{equation}\label{eq:chebyshev}
 |t|=\frac2\pi+\frac4\pi\sum_{k\ge1}
       \frac{(-1)^{k+1}T_{2k}(t)}{4k^2-1}.
\end{equation}
One may obtain this by substituting $t=\cos\theta$ and computing the
cosine coefficients of $|\cos\theta|$. The tail after $k=h$ is bounded
uniformly by
\[
 \frac4\pi\sum_{k>h}\frac1{4k^2-1}
 =\frac2{\pi(2h+1)},
\]
using the telescoping identity
$(4k^2-1)^{-1}=\tfrac12((2k-1)^{-1}-(2k+1)^{-1})$.
The truncation of \eqref{eq:chebyshev} is an even polynomial, hence a
polynomial of degree $h$ in $t^2$. Apply it to $|F_C|/m$ and to
$|F_{\phi(C)}|/m$. The integrals of all polynomial terms agree, and
the two approximation errors together are at most
$4/[\pi(2h+1)]$. Multiplying by $m$ proves \eqref{eq:transport}.
The energy equality is the second moment identity above. For sumsets,
map $a+b$ to $\phi(a)+\phi(b)$: the order-two equivalence makes this
map well-defined and bijective, proving the other assertion.
\end{proof}

\begin{proof}[Completion of Theorem~\ref{thm:inverse} in a high-order model]
Suppose $|A|=N\ge2$ and $\phi:A\to\widetilde A$ is an order-$N$
Freiman isomorphism with $\widetilde A\in\mathfrak G_{\beta,\Delta}$.
Lemma~\ref{lem:transport} gives
\[
 \La(\widetilde A)\le\La(A)+2/\pi
                  \le(K+1)\log N.
\]
Apply the already established theorem to $\widetilde A$ with parameter
$K+1$. Pulling the pieces back by $\phi$ preserves their cardinalities,
disjointness, covering error, and doubling constants by
\eqref{eq:exact-transport}. This proves the stated extension.
\end{proof}

\begin{remark}[The budget is not charged once per model change]
If successive models all belong to the same order-$N$ relation class,
Lemma~\ref{lem:transport} compares any two stages directly. The uniform
error for every corresponding subset is less than $2/\pi$, regardless
of the number of changes. This is the relevant difference from applying
a rough norm-distortion estimate separately at each step. It does not
assert that a geometric assembly can always be found in that relation
class. Minimal-model arguments, such as those developed by
Bedert~\cite{Bedert}, address the size of an integer representative;
that is a separate issue from membership in
$\mathfrak G_{\beta,\Delta}$.
\end{remark}

\begin{corollary}[Small exceptional perturbations]\label{cor:perturbations}
Fix $H_0\ge0$. Suppose $|A|=N$, $\La(A)\le K\log N$, and there is
$A_0\in\mathfrak G_{\beta,\Delta}$, or a set with a corresponding
order-$|A_0|$ model in that class, such that
\[
 |A\mathbin\triangle A_0|\le H_0(\log N)^2.
\]
Then the conclusion of Theorem~\ref{thm:inverse} holds for $A$, with
constants also allowed to depend on $H_0$.
\end{corollary}
\begin{proof}
Put $s=|A\triangle A_0|$. The nonzero Fourier coefficients of
$F_A-F_{A_0}$ are signs on the symmetric difference, so Parseval gives
\[
 |\La(A)-\La(A_0)|\le\norm{F_A-F_{A_0}}_1\le\sqrt s.
\]
Thus $|A_0|=N+o(N)$ and
$\La(A_0)\le K'\log|A_0|$ for a fixed
$K'=O(1+K+\sqrt{H_0})$. Apply Theorem~\ref{thm:inverse} to $A_0$,
using, for example, error $\eps/4$. Intersect each resulting piece
with $A$. At most $s=o(N)$ points are lost in total. Since the number
of pieces is bounded and each originally has size at least $cN$,
each intersection has at least half its former size for sufficiently
large $N$. Its sumset cannot increase, so its doubling constant
increases by at most a factor two. Points of $A\setminus A_0$, along
with the original exceptional set, fit within the prescribed $\eps N$
budget after increasing the size threshold.
\end{proof}

\Needspace{7\baselineskip}
\section{Examples, parameter obstructions, and the general inverse problem}
\label{sec:scope}

\subsection{Interlaced progressions and the meaning of the energy exponent}

The following example explains why the structural conclusion should
not be formulated only through energetic initial segments. It is also
a test that the assembly hypothesis alone is not already the conclusion.

\begin{proposition}[Sidon fibres]\label{prop:sidon-fibres}
For all positive integers $R,M$, there is an integer set $A$ of size
$N=RM$ with
\begin{equation}\label{eq:sidon-fibre-norm}
 \La(A)\le\sqrt R\,\La(\{0,\ldots,M-1\}),
\end{equation}
and, for every $Y\subseteq A$,
\begin{equation}\label{eq:sidon-fibre-energy}
 E(Y)\le2M|Y|^2.
\end{equation}
Every initial segment $P\subseteq A$ with $n=|P|\ge R$ also satisfies
\begin{equation}\label{eq:prefix-obstruction}
 E(P)\le16n^3/R.
\end{equation}
The set belongs to $\mathfrak G_{\beta,\Delta}$ for every fixed
$0<\beta\le1$, $\Delta\ge1$.
\end{proposition}
\begin{proof}
Take $B=\{1,3,\ldots,3^{R-1}\}$, $q=3^R$, and
$A=B+q\{0,\ldots,M-1\}$. The base-three expansion shows that $B$
is Sidon as an integer set; since $q>2\max B$, it is also Sidon
modulo $q$. In particular the displayed parametrization is disjoint.
Writing $D_M(x)=\sum_{j=0}^{M-1}e(jx)$, finite Parseval gives, for
fixed $x$,
\[
 \frac1q\sum_{j=0}^{q-1}|F_B((j+x)/q)|
 \le\left(\frac1q\sum_{j=0}^{q-1}|F_B((j+x)/q)|^2\right)^{1/2}
 =\sqrt R.
\]
Since $F_A(\theta)=F_B(\theta)D_M(q\theta)$, integration proves
\eqref{eq:sidon-fibre-norm}.

Write $Y=\bigcup_{b\in B}(b+qY_b)$, where
$Y_b\subseteq\{0,\ldots,M-1\}$. In any additive quadruple from $Y$,
the residue labels on the two sides pair trivially. For two fibres,
\[
 E(Y_b,Y_c):=\int|F_{Y_b}|^2|F_{Y_c}|^2
 \le\min\{|Y_b|,|Y_c|\}|Y_b||Y_c|
 \le M|Y_b||Y_c|.
\]
For the first inequality, use the supremum bound for one Fourier
factor and Parseval for the other, and choose the smaller result.
Summing over the two possible pairings of residue labels proves
\eqref{eq:sidon-fibre-energy}; the double count when all labels agree
only makes an upper bound larger.

For an initial segment of size $n\ge R$, put $m=\lfloor n/R\rfloor$.
All its points lie in $B+q\{0,\ldots,m\}$. There are at most $2R^2$
ordered residue quadruples, and at most $(m+1)^3$ vertical solutions
for each, because three vertical coordinates determine the fourth.
As $m+1\le2n/R$, this gives \eqref{eq:prefix-obstruction}.
Finally, the root has Sidon residue support $B$, and each child is an
interval. Intervals are long-block terminals with $V=1$, so $A$ belongs
to the stated class.
\end{proof}

Taking $K=C\sqrt R$ in this example and letting $M$ tend to infinity
shows that no universal initial-segment statement can replace the
$K^{-2}$ scale by $K^{-p}$ for a fixed $p<2$, with a positive coefficient
independent of $K$. For fixed $\delta>0$, choose $M$ sufficiently large
that every initial segment of size $\gg N^{1-\delta}$ has size at least
$R$, and then let $R\to\infty$ in \eqref{eq:prefix-obstruction}.
This is a statement about the parameter dependence of that statistic,
not about extremizers for the sharp Littlewood constant.

Indeed each individual fibre is an arithmetic progression of length
$M=N/R$, with energy $(2M^3+M)/3$. For fixed $K$, selecting such a
fibre already gives a linear-sized structured piece in this example.
Conversely, \eqref{eq:sidon-fibre-energy} shows that as $R\to\infty$
the class membership alone cannot give uniform positive-density
small-doubling pieces: if $|Y|\ge cN$ and $|Y+Y|\le Q|Y|$, then
\eqref{eq:energy-doubling} forces $Q\ge cR/2$. The low-$L^1$ assumption
is therefore essential to Theorem~\ref{thm:inverse}.

There is also a parameter restriction for statements that allow
arbitrary subsets. The same example gives
$|Y|\omega(Y)\le2N/R$. If both $K$ and $R$ grow with $N$, any proposed
improvement must account for the corresponding size cost. This does
not preclude a fixed-$K$ linear-sized inverse theorem, which is the
structural direction considered here.

\subsection{Depth is genuinely unbounded}

Let $h,M\ge1$ and consider
\begin{equation}\label{eq:unbalanced-example}
 A=3^h\{0,\ldots,M-1\}\ \cup\ \{3^j:0\le j<h\}.
\end{equation}
Successively reducing modulo $3$ gives a chain of $h$ Sidon splits
with residue support $\{0,1\}\subset\Z/3\Z$. At each step there is
one continuing child and one singleton. The terminal continuing set
is an interval. On the other hand, the triangle inequality and Parseval
give
\[
 \La(A)\le\La(\{0,\ldots,M-1\})+\sqrt h.
\]
For $h=\lfloor(\log M)^2\rfloor$, this is $O(\log|A|)$ while the
depth tends to infinity. A per-step multiplicative loss would not
suffice for this family. Proposition~\ref{prop:leaf-budget} has no such
loss.

For comparison, the fully branching ternary digit set
\[
 \left\{\sum_{j=0}^{h-1}\epsilon_j3^j:\epsilon_j\in\{0,1\}\right\}
\]
has $2^h$ singleton leaves in a Sidon tree. The leaf budget implies
$\La(A)\ge|A|^{1/20}$, incompatible with $K\log|A|$ for large $|A|$
and fixed $K$. Thus the same estimate distinguishes many macroscopic
independent branches from a long chain with only small exceptional
branches.

\subsection{What the proof does not supply}

The missing step for an unrestricted inverse theorem is not part of
BSG, the majority calculation, or the leaf budget. It is a theorem
placing arbitrary low-$L^1$ sets, or suitable high-order equivalent
models, into a usable assembly. Non-Sidon quotient supports need not
be dense separated blocks or long arithmetic inflations. The absence
of a Sidon split is not a classification of a general integer set.

A related difficulty is the difference between localizing a test
function and localizing the target. The kernel $V_D$ in
Lemma~\ref{lem:band} has $2D+1$ Fourier coefficients equal to one but
$\norm{V_D}_1\le3$, because its other coefficients form fractional
slopes. Thus retaining many unit coefficients in a smooth projection
does not justify replacing that projection by an integer indicator
with comparable norm. All target sets in the local theorems above
are actual integer sets, and all residue supports at internal vertices
are the full supports. These requirements are used explicitly.

A natural unrestricted target would assert
\eqref{eq:cover-conclusion} for every $A$ with
$\La(A)\le K\log N$. It would imply a linear-sized small-doubling
subset and positive normalized energy for every macroscopic subset.
The signed-progression conjecture formulated in \cite[Section 6]{BG}
asks for a still more specific representation, with boundedly many
progressions and a sublinear exceptional set. Neither general statement
is proved here. The present theorem isolates local mechanisms that
are compatible with arbitrarily deep independent assemblies, and
makes their exact entry hypotheses explicit.

\appendix

\Needspace{7\baselineskip}
\section{An exact certificate for the numerical coefficient}
\label{app:certificate}

This appendix specifies a finite rational computation certifying the
decimal in Theorem~\ref{thm:constant}. Decimal evaluations are used only
for presentation. Every inequality required by the certificate is an
inequality between rational numbers, or follows from squaring two
nonnegative rational bounds.

The even moments $I_{2r}$ are given exactly by
\eqref{eq:even-sinc-formula}. For the odd coefficients, choose an upper
rational square-root bound as follows. If $x\ge0$ is rational, let
\[
 U_d(x)=\frac{\left\lfloor10^d\sqrt{x}\right\rfloor+1}{10^d}.
\]
The integer in this definition is computed by integer square root after
clearing denominators. Then $\sqrt x<U_d(x)$, with an error at most
$10^{-d}$. Set $\overline c_{2r}=I_{2r}$ and
$\overline c_{2r+1}=U_{40}(I_{2r}I_{2r+2})$.
Using $c_k\le1$, for any $L\ge5$ one has
\begin{equation}\label{eq:series-certificate}
 \sum_{k\ge5}\frac{c_ka^k}{k}
 \le\sum_{k=5}^L\frac{\overline c_ka^k}{k}
        +\frac{a^{L+1}}{(L+1)(1-a)}.
\end{equation}
We use $L=160$ and $a=9/10$.

Logarithms are bounded by their elementary positive series. For a
rational $x>1$, write $z=(x-1)/(x+1)$ and
\[
 \ell_J(x)=2\sum_{j=0}^{J-1}\frac{z^{2j+1}}{2j+1}.
\]
Then
\begin{equation}\label{eq:log-certificate}
 \ell_J(x)<\log x
 <\ell_J(x)+\frac{2z^{2J+1}}{(2J+1)(1-z^2)}.
\end{equation}
The certificate takes $J=240$ for $x=9,10$. These estimates follow by
integrating the geometric series for $(1-z^2)^{-1}$ and bounding each
tail denominator below by $2J+1$.

Substituting \eqref{eq:series-certificate} into \eqref{eq:E-intro},
with all square roots rounded upwards, gives a rational number
$\overline{\mathcal E}$ satisfying
\begin{equation}\label{eq:E-certified}
 \mathcal E_{9/10}\le\overline{\mathcal E}
 <\frac{479084457020255861}{10^{18}}.
\end{equation}
The series remainder in \eqref{eq:series-certificate} is less than
$2.668198549256\cdot10^{-9}$. Its contribution to
$\overline{\mathcal E}$ is multiplied by $2/9$.
Let $\overline L_9$ be the upper bound for $\log9$ in
\eqref{eq:log-certificate}, and let
$\overline R=U_{40}(\overline{\mathcal E}/3)$. Then
\begin{equation}\label{eq:c-certified}
 c_*\ge\frac{9}{10\overline L_9}(1-\overline R)
 >\frac{2459209213577253}{10^{16}}>0.2459209.
\end{equation}
The last inequality is checked by integer cross-multiplication; its
numerator $1-\overline R$ is positive. The same logarithmic enclosure
verifies that \eqref{eq:Phi-endpoints} at $b=11/8$ is negative, with
an upper bound less than $-1.1121923140$.

The companion file \texttt{verify\_constant.py} implements these steps
using integer arithmetic, rational fractions, and integer square roots.
Its printed decimal strings are not inputs to the certificate. The
source and the certified rational procedure are supplied to make the
explicit constant reproducible, rather than to replace the analytical
argument preceding this appendix.

\section{The one-sided endpoint barrier and a two-sided interval model}
\label{app:landau}

This appendix distinguishes the sharp Littlewood constant from the
best value obtainable in the final spectral class of the outer
recursion. The upper bound below is the classical Landau inequality;
see \cite{LandauAsymptotics,Kovalev}. We include its short proof to make
the limitation of the method precise.

\begin{proposition}[Landau's bound in the relevant spectral class]
\label{prop:landau}
Let $D_N(z)=1+z+\cdots+z^{N-1}$ on the unit circle. If
$\norm G_\infty\le1$ and $\supp\widehat G\subseteq(-\infty,N-1]\cap\Z$,
then
\begin{equation}\label{eq:landau}
 |\ip{D_N}G|\le
 \sum_{k=0}^{N-1}\left(\frac{\binom{2k}{k}}{4^k}\right)^2
 =\frac1\pi\log N+O(1).
\end{equation}
The inequality is sharp within this class of tests.
\end{proposition}
\begin{proof}
Put $m=N-1$, $b_k=\binom{2k}{k}/4^k$, and
$Q_m(z)=\sum_{k=0}^m b_kz^k$. Since the $b_k$ are the coefficients of
$(1-z)^{-1/2}$,
\[
 Q_m(z)^2=1+z+\cdots+z^m+O(z^{m+1}).
\]
The boundary function $h=z^m\overline G$ is analytic and bounded by
one in the disc. If $h_k$ are its Taylor coefficients, then
\[
 \ip{D_N}G=\sum_{k=0}^m h_k=[z^m](hQ_m^2).
\]
Taking this coefficient on the unit circle gives
\[
 |[z^m](hQ_m^2)|\le\int_\T|Q_m|^2=\sum_{k=0}^m b_k^2.
\]
Stirling's formula gives $b_k^2=1/(\pi k)+O(k^{-2})$ for $k\ge1$,
proving the asymptotic.

For sharpness, $Q_m$ has no zeros in the closed disc. This is immediate
for $m=0$. For $m\ge1$, the identity
\[
 (1-z)Q_m(z)=1-\sum_{k=1}^m(b_{k-1}-b_k)z^k-b_mz^{m+1}
\]
has positive subtracted coefficients summing to one. Their sum has
modulus strictly less than one when $|z|<1$. Equality on the unit
circle can produce a zero of the right side only at $z=1$, since the
coefficient of $z$ is positive; but $Q_m(1)>0$. Thus
$h=z^m\overline{Q_m}/Q_m$, with the numerator interpreted as the
reversed polynomial, is analytic and has unit boundary modulus. It
attains the bound, since
$hQ_m^2=z^m|Q_m|^2$ on the circle. The corresponding test is
$G=Q_m/\overline{Q_m}$.
\end{proof}

For an interval, every test produced in \eqref{eq:test-function} has
Fourier support at most $N-1$. Consequently modifying source phases,
scale choices, or multiplier estimates while retaining that final
support restriction cannot prove a coefficient greater than $1/\pi$.
This is not an upper bound for the actual Littlewood constant, and it
does not prohibit using one-sided tests as inputs to a nonlinear
construction with two-sided spectrum.

\begin{proposition}[Two-sided completion on an interval]
\label{prop:interval-completion}
With $Q_{N-1}$ as above, put
\[
 U_N=\frac{Q_{N-1}}{\overline{Q_{N-1}}},\qquad
 V_N=z^{N-1}\overline{U_N},\qquad
 H_N=\frac{U_N+V_N}{|U_N+V_N|},
\]
where $H_N=0$ at zeros of its denominator. Then $|H_N|\le1$ and
\begin{equation}\label{eq:interval-completion}
 \Rea\ip{D_N}{H_N}\ge\norm{D_N}_1-O(1)
 =\frac4{\pi^2}\log N+O(1).
\end{equation}
\end{proposition}
\begin{proof}
In this proof write $z=e^{it}$ with $0<t<2\pi$; circle integrals carry
measure $dt/(2\pi)$. Define
\[
 U_\infty(t)=ie^{-it/2},\qquad
 V_\infty(t)=-ie^{i(N-1/2)t}.
\]
The identity
\begin{equation}\label{eq:interval-phase-identity}
 U_\infty(t)+V_\infty(t)=2\sin(t/2)D_N(e^{it})
\end{equation}
shows that the phase of this sum equals the phase of $D_N$. These are
boundary phase models, not assertions about their one-sided Fourier
support.

Abel summation applied to the decreasing coefficients $b_k$ gives,
on the unit circle away from $1$,
\[
 |Q_{N-1}(z)-(1-z)^{-1/2}|
 \le C N^{-1/2}|1-z|^{-1}.
\]
For example the tail of $\sum_{k\ge N}b_kz^k$ is bounded by a
constant times $b_N/|1-z|$ by summation by parts. Since
$|(1-z)^{-1/2}|=|1-z|^{-1/2}$, normalization of complex numbers implies,
with $s=\min(t,2\pi-t)$,
\begin{equation}\label{eq:phase-error}
 |U_N-U_\infty|+|V_N-V_\infty|
 \le C\min\{1,(Ns)^{-1/2}\}.
\end{equation}
Here the phase of $(1-e^{it})^{-1/2}$ divided by its conjugate is
exactly $ie^{-it/2}$.

On the arcs $s\le1/N$, use $|D_N|\le N$ and the trivial bound two
for the difference of normalized phases. The weighted error integral
is $O(1)$. Outside these arcs, use
\[
 \left|\frac X{|X|}-\frac Y{|Y|}\right|
 \le\frac{2|X-Y|}{|Y|},\qquad
 \frac{|D_N|}{|U_\infty+V_\infty|}
 =\frac1{2\sin(t/2)}.
\]
At zeros the inequality is interpreted by an arbitrary unit-bounded
value, and the exceptional finite set has measure zero. Combining
with \eqref{eq:phase-error}, the remaining weighted error is at most
\[
 C N^{-1/2}\int_{1/N}^\pi s^{-3/2}\,ds=O(1).
\]
This proves the first inequality in \eqref{eq:interval-completion};
the second is \eqref{eq:ap-asymptotic}.
\end{proof}

For a general frequency set, nonlinear phase completion introduces
mixed products such as $U^2\overline V$ and $V^2\overline U$ in the
pairing. Opposite one-sided Fourier supports alone do not eliminate
these terms or control their signs. Proposition~\ref{prop:interval-completion}
is an interval model, not a proof of the sharp coefficient for arbitrary
sets.

\section*{Computational reproducibility}
The explicit decimal is certified by the rational procedure in
Appendix~\ref{app:certificate}; accompanying finite checks of the
moment and Sidon identities are supplementary tests, not substitutes
for the proofs.

\section*{Statement on the use of AI}
Generative AI tools were used substantially during exploratory work and
in preparing the manuscript, including the development and checking of
arguments, calculations, exposition, and \LaTeX{} and code preparation.
The author is responsible for checking every argument and for the
mathematical content of any submitted version.

\begin{thebibliography}{99}

\bibitem{Bedert}
B. Bedert, \emph{Large sum-free subsets of sets of integers via
$L^1$-estimates for trigonometric series}, preprint (2025),
\arxiv{2502.08624}.

\bibitem{BG}
T. F. Bloom and B. Green, \emph{Remarks on the inverse Littlewood
conjecture}, Q. J. Math. (2026), article haag016,
\doi{10.1093/qmath/haag016}; \arxiv{2602.16482v2}
(revised April 18, 2026).

\bibitem{Bourgain}
J. Bourgain, \emph{Estimates related to sumfree subsets of sets of
integers}, Israel J. Math. \textbf{97} (1997), 71--92.

\bibitem{BTK}
N. G. de Bruijn, C. van Ebbenhorst Tengbergen, and D. Kruyswijk,
\emph{On the set of divisors of a number}, Nieuw Arch. Wisk. (2)
\textbf{23} (1951), 191--193.

\bibitem{Duren}
P. L. Duren, \emph{Theory of $H^p$ Spaces}, Pure and Applied Mathematics,
vol.~38, Academic Press, New York--London, 1970.

\bibitem{Fournier}
J. J. F. Fournier, \emph{Some remarks on the recent proofs of the
Littlewood conjecture}, in \emph{Second Edmonton Conference on
Approximation Theory} (Edmonton, Alta., 1982), CMS Conf. Proc.,
vol.~3, American Mathematical Society, Providence, RI, 1983,
157--170.

\bibitem{Gabriel}
R. M. Gabriel, \emph{The rearrangement of positive Fourier coefficients},
Proc. London Math. Soc. (2) \textbf{33} (1932), 32--51,
\doi{10.1112/plms/s2-33.1.32}.

\bibitem{GreenSanders}
B. Green and T. Sanders, \emph{A quantitative version of the idempotent
theorem in harmonic analysis}, Ann. of Math. (2) \textbf{168} (2008),
no.~3, 1025--1054, \doi{10.4007/annals.2008.168.1025}.

\bibitem{Hanson}
B. Hanson, \emph{Littlewood's problem for sets with multidimensional
structure}, Int. Math. Res. Not. IMRN \textbf{2021} (2021), no.~21,
16736--16750, \doi{10.1093/imrn/rnaa043}; \arxiv{2003.01561}.

\bibitem{HardyLittlewood}
G. H. Hardy and J. E. Littlewood, \emph{A new proof of a theorem on
rearrangements}, J. London Math. Soc. \textbf{23} (1948), 163--168.

\bibitem{Konyagin}
S. V. Konyagin, \emph{On the Littlewood problem}, Izv. Akad. Nauk SSSR
Ser. Mat. \textbf{45} (1981), no.~2, 243--265, 463; English translation,
Math. USSR-Izv. \textbf{18} (1982), no.~2, 205--225.

\bibitem{Kovalev}
L. V. Kovalev, \emph{Sharp bounds for some segments of bounded power
series}, preprint (2025), \arxiv{2507.04544}.

\bibitem{LandauAsymptotics}
Y. Li, S. Liu, S. Xu, and Y. Zhao, \emph{Asymptotics of Landau constants
with optimal error bounds}, Constr. Approx. \textbf{40} (2014),
281--305, \doi{10.1007/s00365-014-9259-x}.

\bibitem{MPS}
O. C. McGehee, L. Pigno, and B. Smith, \emph{Hardy's inequality and the
$L^1$ norm of exponential sums}, Ann. of Math. (2) \textbf{113} (1981),
no.~3, 613--618, \doi{10.2307/2007000}.

\bibitem{Petridis}
G. Petridis, \emph{The $L^1$-norm of exponential sums in $\Z^d$},
Math. Proc. Cambridge Philos. Soc. \textbf{154} (2013), no.~3,
381--392, \doi{10.1017/S0305004112000588}.

\bibitem{Pichorides}
S. K. Pichorides, \emph{On the $L^1$ norm of exponential sums},
Ann. Inst. Fourier (Grenoble) \textbf{30} (1980), no.~2, 79--89,
\doi{10.5802/aif.785}.

\bibitem{ReiherSchoen}
C. Reiher and T. Schoen, \emph{Note on the theorem of Balog,
Szemer\'edi, and Gowers}, Combinatorica \textbf{44} (2024), no.~3,
691--698, \doi{10.1007/s00493-024-00092-5}; \arxiv{2308.10245v2}.

\bibitem{Stegeman}
J. D. Stegeman, \emph{On the constant in the Littlewood problem},
Math. Ann. \textbf{261} (1982), no.~1, 51--54,
\doi{10.1007/BF01456409}.

\bibitem{TaoVu}
T. Tao and V. H. Vu, \emph{Additive Combinatorics}, Cambridge Studies
in Advanced Mathematics, vol.~105, Cambridge University Press,
Cambridge, 2006, \doi{10.1017/CBO9780511755149}.

\bibitem{Yabuta}
K. Yabuta, \emph{A remark on the Littlewood conjecture}, Bull. Fac.
Sci. Ibaraki Univ. Ser. A \textbf{14} (1982), 19--21,
\doi{10.5036/bfsiu1968.14.19}.

\bibitem{Zygmund}
A. Zygmund, \emph{Trigonometric Series}, Vols.~I and II, third edition,
Cambridge Mathematical Library, Cambridge University Press,
Cambridge, 2002.

\end{thebibliography}
\end{document}
